\subsection{Multiplication endomorphisms}

Let X be a locally compact topological space and let \(\mu\) be a positive measure on X. Let \(p \in [1, +\infty[\) .

Let \(g \in \mathcal{L}^{\infty}(\mathrm{X},\mu)\) . Denote by \(m_{g}\) the map \(f \mapsto g \cdot f\) of \(\mathcal{L}^{p}(\mathrm{X},\mu)\) into itself. The map \(m_{g}\) is linear and continuous since

\[
\mathrm{N} _ {p} (m _ {g} (f)) \leqslant \mathrm{N} _ {\infty} (g) \mathrm{N} _ {p} (f)\tag{1}
\]

for every function \(f \in \mathcal{L}^{p}(\mathrm{X}, \mu)\) . In particular, the function \(m_{g}(f)\) is \(\mu\) -negligible if f is \(\mu\) -negligible. The map \(m_{g}\) therefore induces, by passing to quotients, an endomorphism of \(\mathrm{L}^{p}(\mathrm{X}, \mu)\) , which will be denoted by \(\widetilde{m}_{g}\) .

Lemma 7. --- Let \(g \in \mathcal{L}^{\infty}(\mathrm{X},\mu)\) . The endomorphism \(\tilde{m}_g\) of \(\mathrm{L}^{p}(\mathrm{X},\mu)\) is injective if and only if the set of \(x \in \mathrm{X}\) such that \(g(x) = 0\) is locally \(\mu\) -negligible.

Let A denote the set of \(x \in X\) such that \(g(x) = 0\) .

Suppose that A is locally \(\mu\) -negligible. Let \(f \in \mathcal{L}^p(\mathrm{X}, \mu)\) and \(\widetilde{f}\) its class in \(\mathrm{L}^p(\mathrm{X}, \mu)\) . Suppose that \(\widetilde{m}_g(\widetilde{f}) = 0\) . By definition, this is the case if and only if the function \(fg\) is \(\mu\) -negligible, so that the set B of \(x \in \mathrm{X}\) such that \(f(x)g(x) \neq 0\) is \(\mu\) -negligible. The function \(f\) vanishes outside \(\mathrm{A} \cup \mathrm{B}\) , hence is \(\mu\) -negligible. This implies that the endomorphism \(\widetilde{m}_g\) is injective.

Conversely, suppose that A is not locally \(\mu\) -negligible. Let C be a compact subset of X such that A∩C is not \(\mu\) -negligible, and let \(\varphi\) be the characteristic function of A∩C. The class of the function \(\varphi\) in L \(^{p}\) (X, \(\mu\) ) is not zero, but that of \(m_{g}(\varphi)\) is, so \(m_{g}\) is not injective.

Since the product of a locally \(\mu\) -negligible function and a \(\mu\) -negligible function is \(\mu\) -negligible, the endomorphism \(\widetilde{m}_g\) depends only on the class of \(g\) in \(\mathrm{L}^{\infty}(\mathrm{X},\mu)\) . For \(\widetilde{g}\in \mathrm{L}^{\infty}(\mathrm{X},\mu)\) , one also denotes by \(\widetilde{m}_{\widetilde{g}}\) the endomorphism \(\widetilde{m}_g\) of \(\mathrm{L}^p (\mathrm{X},\mu)\) for every function \(g\in \mathcal{L}^{\infty}(\mathrm{X},\mu)\) of which \(\widetilde{g}\) is the class. One says that \(\widetilde{m}_{\widetilde{g}}\) is the multiplication endomorphism by \(\widetilde{g}\) in \(\mathrm{L}^p (\mathrm{X},\mu)\) .

PROPOSITION 5. --- Let \(p \in [1, +\infty[\) . The map \(m: g \mapsto \widetilde{m}_g\) from \(\mathrm{L}^{\infty}(\mathrm{X}, \mu)\) into \(\mathcal{L}(\mathrm{L}^{p}(\mathrm{X}, \mu))\) is an isometric morphism of unital Banach algebras.

For \(a\) , \(b\) in \(\mathbf{C}\) and \(g_{1}\) , \(g_{2}\) in \(\mathrm{L}^{\infty}(\mathrm{X},\mu)\) , one immediately verifies that \(\widetilde{m}_{ag_1+bg_2}=a\widetilde{m}_{g_1}+b\widetilde{m}_{g_2}\) , so the map \(\widetilde{m}\) is linear. Moreover, one has \(\widetilde{m}_{1}=1_{\mathrm{L}^{p}(\mathrm{X},\mu)}\) and \(\widetilde{m}_{g_{1}g_{2}}=\widetilde{m}_{g_{1}}\widetilde{m}_{g_{2}}\) , so the map \(\widetilde{m}\) is a unital algebra morphism from \(\mathrm{L}^{\infty}(\mathrm{X},\mu)\) into \(\mathcal{L}(\mathrm{L}^{p}(\mathrm{X},\mu))\) .

Formula (1) above shows that \(\widetilde{m}\) has norm \(\leqslant 1\) .

Let \(g \in \mathcal{L}^{\infty}(\mathrm{X},\mu)\) and \(\widetilde{g}\) be its class in \(\mathrm{L}^{\infty}(\mathrm{X},\mu)\) . Let \(\varepsilon > 0\) . The set Y of \(x \in X\) such that \(|g(x)| > \| \widetilde{g}\|_{\infty} - \varepsilon\) is not locally \(\mu\) -negligible. There therefore exists a compact subset C of X such that \(\mu (\mathrm{Y} \cap \mathrm{C}) > 0\) . Let \(\varphi\) be the characteristic function of \(\mathrm{Y} \cap \mathrm{C}\) . Since

\[
\begin{array}{c} \| \widetilde {m} _ {\widetilde {g}} (\varphi) \| _ {p} = \Bigl (\int_ {\mathrm{X}} | \varphi g | ^ {p}   d \mu \Bigr) ^ {1 / p} \geqslant \Bigl (\int_ {\mathrm{Y}} | \varphi g | ^ {p}   d \mu \Bigr) ^ {1 / p} \\ \geqslant (\| \widetilde {g} \| _ {\infty} - \varepsilon) \mu (\mathrm{Y} \cap \mathrm{C}) ^ {1 / p} = (\| \widetilde {g} \| _ {\infty} - \varepsilon) \| \varphi \| _ {p}, \end{array}
\]

it follows that \(\|\widetilde{m}_{\widetilde{g}}\|\geqslant\|\widetilde{g}\|_{\infty}-\varepsilon\) . Since \(\varepsilon\) is arbitrary, we deduce that \(\|\widetilde{m}_{\widetilde{g}}\|\geqslant\|\widetilde{g}\|_{\infty}\) , which proves that \(\widetilde{m}\) is isometric.

Lemma 8. --- Let \((g_n)\) be a bounded sequence in \(\mathcal{L}^\infty(\mathrm{X},\mu)\) converging simply \(\mu\) -almost everywhere. Let \(\widetilde{m} \in \mathrm{L}^\infty(\mathrm{X},\mu)\) be the class of its limit. Then \(\widetilde{m}_{g_n}\) converges to \(\widetilde{m}_{\widetilde{g}}\) in the space \(\mathcal{L}(\mathrm{L}^p(\mathrm{X},\mu))\) endowed with the topology of simple convergence.

Let \(f \in \mathcal{L}^p(\mathrm{X},\mu)\) . The sequence \((g_{n}f)\) is bounded in \(\mathcal{L}^p(\mathrm{X},\mu)\) and converges simply \(\mu\) -almost everywhere to \(gf\) . By the Lebesgue theorem (INT, IV, p. 137, § 3, no. 7, th. 6), the sequence \((g_{n}f)\) converges to \(gf\) in \(\mathcal{L}^p(\mathrm{X},\mu)\) , and the assertion follows.

We now consider the case p = 2.

PROPOSITION 6. --- The map m: \(g \mapsto \widetilde{m}_{g}\) is a unital isometric morphism of the stellar algebra \(\mathrm{L}^{\infty}(\mathrm{X}, \mu)\) into the stellar algebra \(\mathcal{L}(\mathrm{L}^{2}(\mathrm{X}, \mu))\) .

In particular, for every \(g \in \mathrm{L}^{\infty}(\mathrm{X}, \mu)\) the multiplication endomorphism \(\widetilde{m}_{g}\) is normal; it is Hermitian if and only if g is locally real-valued \(\mu\) -almost everywhere.

By prop. 5, the map \(\tilde{m}\) is an injective, isometric morphism of unital Banach algebras from \(\mathrm{L}^{\infty}(\mathrm{X},\mu)\) into \(\mathcal{L}(\mathrm{L}^{2}(\mathrm{X},\mu))\) . Let \(g\in \mathcal{L}^{\infty}(\mathrm{X},\mu)\) . For \(f_{1}\) and \(f_{2}\in \mathcal{L}^{2}(\mathrm{X},\mu)\) , we have

\[
\langle f _ {1} \mid \widetilde {m} _ {g} (f _ {2}) \rangle = \int_ {\mathrm{X}} \overline {{f _ {1} (x)}} g (x) f _ {2} (x) d \mu (x) = \langle \widetilde {m} _ {\overline {{g}}} (f _ {1}) \mid f _ {2} \rangle ,
\]

from which it follows that \(\widetilde{m}_g^* = \widetilde{m}_{\overline{g}}\) , which proves that \(m\) is an involutive morphism. The last assertions follow from this (cf. I, p. 106, prop. 5).

COROLLARY. --- Let \(g \in \mathrm{L}^{\infty}(\mathrm{X},\mu)\) .

\begin{enumerate}
\def\labelenumi{\alph{enumi})}
\item
  The spectrum of \(\widetilde{m}_g\) in \(\mathcal{L}(\mathrm{L}^2 (\mathrm{X},\mu))\) is the \(\mu\) -essential image of \(g\) ;
\item
  The endomorphism \(\widetilde{m}_g\) is positive if and only if \(g\) is locally positive \(\mu\) -almost everywhere;
\item
  For any function \(f \in \mathcal{C}(\mathrm{Sp}(\widetilde{m}_g))\) , we have \(f(\widetilde{m}_g) = \widetilde{m}_{f \circ g}\) .
\end{enumerate}

Since \(\widetilde{m}\) is injective, we have \(\mathrm{Sp}(\widetilde{m}_{g}) = \mathrm{Sp}(g)\) by prop. 5 of I, p. 106; the result then follows from lemma 6 of IV, p. 185, proposition 4 of IV, p. 185, and definition 6 of I, p. 115.

PROPOSITION 7. --- The image of the morphism \(\tilde{m}\) from \(\mathrm{L}^{\infty}(\mathrm{X},\mu)\) into \(\mathcal{L}(\mathrm{L}^{2}(\mathrm{X},\mu))\) is a maximal commutative subalgebra of the algebra \(\mathcal{L}(\mathrm{L}^{2}(\mathrm{X},\mu))\) .

It suffices to prove that if \(u \in \mathcal{L}(\mathrm{L}^{2}(\mathrm{X}, \mu))\) is an endomorphism that commutes with \(\widetilde{m}_{g}\) for every function \(g \in \mathcal{L}^{\infty}(\mathrm{X}, \mu)\) , then u belongs to the image of \(\widetilde{m}\) . Let u be such an endomorphism.

We denote \(\tilde{f}\) the class in \(\mathrm{L}^2 (\mathrm{X},\mu)\) of a function \(f\in \mathcal{L}^2 (\mathrm{X},\mu)\) . Denote by \(v\) the linear map from \(\mathcal{L}^2 (\mathrm{X},\mu)\) into \(\mathrm{L}^2 (\mathrm{X},\mu)\) defined by \(f\mapsto u(\widetilde{f})\) . We have \(v\circ m_g = \widetilde{m}_g\circ v\) for all \(g\in \mathcal{L}^{\infty}(\mathrm{X},\mu)\) .

For any subset \(\mu\) -measurable Y of X, we denote \(\varphi_{\mathrm{Y}}\) its characteristic function; the endomorphism \(\widetilde{m}_{\varphi_{\mathrm{Y}}}\) is an orthogonal projector of \(\mathrm{L}^2 (\mathrm{X},\mu)\) .

Suppose first that X is compact, so that the class of the constant function 1 belongs to \(\mathrm{L}^{2}(\mathrm{X},\mu)\) . Let g be a function in \(\mathcal{L}^{2}(\mathrm{X},\mu)\) whose class in \(\mathrm{L}^{2}(\mathrm{X},\mu)\) is \(v(1)\) .

Let \(c\) be a positive real number and Y the set of \(x \in \mathrm{X}\) such that \(|g(x)| \geqslant c\) ; it is a \(\mu\) -measurable set. We have in \(\mathrm{L}^2(\mathrm{X}, \mu)\) the equalities \(\widetilde{\varphi_{\mathrm{Y}}} g = \widetilde{m}_{\varphi_{\mathrm{Y}}}(v(1)) = v(m_{\varphi_{\mathrm{Y}}}(1)) = v(\varphi_{\mathrm{Y}})\) .

Since moreover \(|c\varphi_{\mathrm{Y}}| \leqslant |g\varphi_{\mathrm{Y}}|\) , we obtain

\[
c ^ {2} \mu (\mathrm{Y}) = \int_ {\mathrm{X}} (c \varphi_ {\mathrm{Y}}) ^ {2} d \mu \leqslant \int_ {\mathrm{X}} | \varphi_ {\mathrm{Y}} g | ^ {2} d \mu = \| \widetilde {u} (\varphi_ {\mathrm{Y}}) \| ^ {2} \leqslant \| u \| ^ {2} \mu (\mathrm{Y}).
\]

This inequality implies that \(\mu(\mathrm{Y}) = 0\) if \(c > \|u\|\) , so that the function \(g\) is bounded \(\mu\) -almost everywhere by \(\|u\|\) . For every function \(f \in \mathcal{C}(\mathrm{X})\) , we finally have

\[
u (\widetilde {f}) = v (m _ {f} (1)) = \widetilde {m} _ {f} (v (1)) = \widetilde {m} _ {f} (\widetilde {g}) = \widetilde {m} _ {g} (\widetilde {f}),
\]

whence the equality \(u = \tilde{m}_{g}\) , and in particular \(\mathrm{N}_{\infty}(g) = \| u\|\) .

Consider now the general case. There exists a locally countable family \((\mathrm{X}_{i})_{i\in\mathrm{I}}\) of pairwise disjoint compact subsets of X such that \(Z = X - \bigcup_{i \in I} X_i\) is locally \(\mu\) -negligible (INT, IV, p.190, § 5, no. 9, prop. 14). Let \(\mu_i\) be the measure induced by \(\mu\) on \(X_i\) (INT, IV, p. 186, § 5, no. 7, def. 4).

By prop. 1 of IV, p. 179, for every \(i \in I\) the image \(\mathrm{E}_i\) of \(\widetilde{m}_{\varphi_{\mathrm{X}_i}}\) is identified with \(\mathrm{L}^2(\mathrm{X}_i, \mu_i)\) by \(f \mapsto f|\mathrm{X}_i\) . For every function \(g \in \mathcal{L}^\infty(\mathrm{X}, \mu)\) the multiplication endomorphism \(\widetilde{m}_g\) commutes with \(\widetilde{m}_{\varphi_{\mathrm{X}_i}}\) , hence leaves stable \(\mathrm{E}_i\) , and the endomorphism of \(\mathrm{E}_i\) deduced from \(\widetilde{m}_g\) by passage to the subspaces coincides with the endomorphism of multiplication by \(g|\mathrm{X}_i\) on \(\mathrm{L}^2(\mathrm{X}_i, \mu_i)\) .

Since \(u\) commutes with \(m_{\varphi_{X_i}}\) , it also leaves stable the subspace \(\mathrm{E}_i\) . Denote by \(u_i\) the endomorphism of \(\mathrm{L}^2 (\mathrm{X}_i,\mu_i)\) deduced from \(u\) by passage to the subspaces.

The map from \(\mathcal{L}^{\infty}(\mathrm{X},\mu)\) into \(\mathcal{L}^{\infty}(\mathrm{X}_{i},\mu_{i})\) defined by \(g\mapsto g|\mathrm{X}_{i}\) being surjective, the hypothesis implies that \(u_{i}\) commutes with \(\widetilde{m}_{g}\) for every function \(g\in \mathcal{L}^{\infty}(\mathrm{X}_{i},\mu_{i})\) . By the first part of the proof, there exists a function \(g_{i}\in \mathcal{L}^{\infty}(\mathrm{X}_{i},\mu_{i})\) such that \(u_{i} = \widetilde{m}_{g_{i}}\) and \(\mathrm{N}_{\infty}(g_{i})\leqslant \| u_{i}\| \leqslant \| u\|\) .

The function g on X which coincides with \(g_{i}\) on \(X_{i}\) and which is zero on Z is bounded and \(\mu\) -measurable (INT, IV, p. 193, § 5, no. 10, prop. 16), hence \(g \in \mathcal{L}^{\infty}(X, \mu)\) . For every \(i \in I\) , the restriction of u to \(E_{i}\) coincides with that of \(\tilde{m}_{g}\) ; consequently, we have \(u = m_{g}\) by prop. 1, c) of IV, p. 179.

COROLLARY. --- Let u be an endomorphism of the Hilbert space \(\mathrm{L}^{2}(\mathrm{X},\mu)\) commuting with \(\widetilde{m}_{g}\) for every function \(g\in\mathcal{K}(\mathrm{X})\) . Then there exists a unique element \(f\in\mathrm{L}^{\infty}(\mathrm{X},\mu)\) such that \(u=\widetilde{m}_{f}\) .

By prop. 7, it suffices to prove that \(u\) commutes with \(\widetilde{m}_{g}\) for all \(g\in\mathcal{L}^{\infty}(\mathrm{X},\mu)\) . Let \(h_{1}\) and \(h_{2}\) be elements of \(\mathcal{L}^{2}(\mathrm{X},\mu)\) with classes \(\widetilde{h}_{1}\) and \(\widetilde{h}_{2}\) in \(\mathrm{L}^{2}(\mathrm{X},\mu)\) . Let \(k_{1}\) (resp. \(k_{2}\) ) be a function in \(\mathcal{L}^{2}(\mathrm{X},\mu)\) whose class is \(u(\widetilde{h}_{1})\) (resp. \(u^{*}(\widetilde{h}_{2})\) ).

Denote by \(h = h_1 \overline{k}_2 - k_1 \overline{h}_2\) ; we have \(h \in \mathcal{L}^1(\mathrm{X}, \mu)\) . Define the measure \(\nu = h \cdot \mu\) on \(\mathrm{X}\) ; it is bounded. For every \(g \in \mathcal{L}^\infty(\mathrm{X}, \mu)\) , we have

\[
\begin{array}{c} \langle \widetilde {h} _ {2} | u (\widetilde {m} _ {g} (\widetilde {h} _ {1})) - \widetilde {m} _ {g} (u (\widetilde {h} _ {1})) \rangle = \langle u ^ {*} (\widetilde {h} _ {2}) | \widetilde {m} _ {g} (\widetilde {h} _ {1}) \rangle - \langle \widetilde {h} _ {2} | \widetilde {m} _ {g} (u (\widetilde {h} _ {1})) \rangle \\ = \int_ {\mathrm{X}} g \cdot h _ {1} \cdot \overline {{k}} _ {2} d \mu - \int_ {\mathrm{X}} g \cdot k _ {1} \cdot \overline {{h}} _ {2} d \mu = \nu (g). \end{array}
\]

We therefore have by hypothesis \(\nu(g) = 0\) for every function \(g \in \mathcal{K}(\mathrm{X})\) , that is, \(\nu = 0\) . Consequently, it follows that

\[
\langle \widetilde {h} _ {2} | u (\widetilde {m} _ {g} (\widetilde {h} _ {1})) - \widetilde {m} _ {g} (u (\widetilde {h} _ {1})) \rangle = \nu (g) = 0
\]

for all \(g \in \mathcal{L}^{\infty}(\mathrm{X},\mu)\) . Since this is the case for all \(h_1\) and \(h_2\) in \(\mathcal{L}^2 (\mathrm{X},\mu)\) , we have \(u\circ \widetilde{m}_g = \widetilde{m}_g\circ u\) .

Remark. --- In what follows, we will often simply denote by \(m_{g}\) the multiplication operator by g on \(\mathrm{L}^{p}(\mathrm{X},\mu)\) .

\subsection{Spectral measures}

In this section, we fix a complex Hilbert space E.

Recall that if A is a commutative unital stellar algebra, we denote by X(A) the compact topological space of its unital characters (cor. 1 of I, p. 29) and \(\mathcal{G}_{\mathrm{A}}\) the Gelfand transform of A (def. 5 of I, p. 7), which is an isometric isomorphism of unital involutive algebras from A onto \(\mathcal{C}(\mathrm{X}(\mathrm{A}))\) (th. 1 of I, p. 108). We will denote by \(\mathcal{H}_{\mathrm{A}}\) the inverse isomorphism.

Lemma 9. --- Let A be a commutative unital stellar subalgebra of \(\mathcal{L}(\mathrm{E})\) . For all \(x\) and \(y\) in E, the map \(\mu\) from \(\mathsf{X}(\mathrm{A})\) into \(\mathbf{C}\) defined by \(\mu(f) = \langle x | \mathcal{H}_{\mathrm{A}}(f)y \rangle\) for all \(f \in \mathcal{C}(\mathsf{X}(\mathrm{A}))\) is a bounded measure on the compact space \(\mathsf{X}(\mathrm{A})\) . It is positive if \(x = y\) . Its norm is \(\leqslant \|x\|\ \|y\|\) , with equality if \(x = y\) .

The map \(\mu\) is linear. For every function \(f\in\mathcal{C}(\mathsf{X}(\mathsf{A}))\) , we have

\[
| \mu (f) | \leqslant \| x \| \| y \| \| \mathcal {H} _ {\mathrm{A}} (f) \| = \| x \| \| y \| \| f \|,
\]

hence \(\mu\) is a bounded measure on \(\mathsf{X}(\mathsf{A})\) of norm \(\leqslant \| x\|\| y\|\) .

Suppose that \(x = y\) . For every positive function \(f \in \mathcal{C}(\mathsf{X}(\mathsf{A}))\) , the element \(\mathcal{H}_{\mathsf{A}}(f)\) is a positive element of A (example 1 of I, p. 115), hence a positive element of \(\mathcal{L}(\mathsf{E})\) , whence \(\mu(f) = \langle x | \mathcal{H}_{\mathsf{A}}(f)(x) \rangle \geqslant 0\) (prop. 8 of I, p. 138). This shows that \(\mu\) is a positive measure. Since \(\mu(1) = \|x\|^2\) the total mass of \(\mu\) is \(\|x\|^2\) (INT, III, p. 58, § 1, no. 8, cor. 2).

DEFINITION 2. --- Let A be a commutative unital stellar subalgebra of \(\mathcal{L}(\mathrm{E})\) . For all \(x\) and \(y\) in E, the linear form \(f \mapsto \langle x | \mathcal{H}_{\mathrm{A}}(f)y \rangle\) on \(\mathcal{C}(\mathrm{X}(\mathrm{A}))\) is called the spectral measure of \((x,y)\) relative to A. If \(x = y\) , we say that it is the spectral measure of \(x\) relative to A.

Remark. --- Let A be a commutative unital stellar subalgebra of \(\mathcal{L}(\mathrm{E})\) . For all \(x\) and \(y\) in E, denote \(\mu_{x,y}\) the spectral measure of \((x,y)\) relative to A. The map defined by \((x,y) \mapsto \mu_{x,y}\) from \(\mathrm{E} \times \mathrm{E}\) into \(\mathcal{M}^1(\mathrm{X}(\mathrm{A}))\) is sesquilinear.

Let \(u\) be a normal endomorphism of E and A the unital stellar subalgebra of \(\mathcal{L}(\mathrm{E})\) generated by \(u\) . It is commutative. The space \(\mathsf{X}(\mathrm{A})\) is identified with \(\mathrm{Sp}(u)\) by the mapping \(\chi \mapsto \chi(u)\) (lemma 10 of I, p. 109). The spectral measure \(\mu_{x,y}\) of \((x,y)\) relative to A is therefore identified with a measure on \(\mathrm{Sp}(u)\) , called the spectral measure of \((x,y)\) relative to \(u\) . For every function \(f \in \mathcal{C}(\mathrm{Sp}(u))\) , we then have

\[
\int_ {\mathrm{Sp} (u)} f d \mu_ {x, y} = \langle x | f (u) y \rangle
\]

according to def. 5 of I, p. 110.

\subsection{Commutative stellar algebras of endomorphisms of a Hilbert space}

DEFINITION 3. --- Let A be an algebra over \(\mathbf{C}\) and E a complex topological vector space. Let π be a representation of A in E. An element x of E is said to be a cyclic vector for π if the set of elements π(a)x for a ∈ A is total in E.

Let \(u \in \mathcal{L}(\mathrm{E})\) be an endomorphism of E. We say that \(x \in \mathrm{E}\) is a cyclic vector for \(u\) if it is a cyclic vector for the identity representation of the stellar subalgebra generated by \(u\) in \(\mathcal{L}(\mathrm{E})\) .

PROPOSITION 8. --- Let E be a complex Hilbert space. Let A be a commutative unital stellar subalgebra of \(\mathcal{L}(\mathrm{E})\) . Let \(x\) be an element of E. Set \(\mathrm{E}_x = \overline{\mathrm{A} \cdot x}\) and let us denote \(\mu_x\) the spectral measure of \(x\) relative to A.

There exists a unique isometric isomorphism \(\theta_x\) of \(\mathrm{L}^2 (\mathsf{X}(\mathrm{A}),\mu_x)\) onto \(\mathrm{E}_x\) such that \(\theta_{x}(f) = \mathcal{H}_{\mathrm{A}}(f)(x)\) for every function \(f\in \mathcal{C}(\mathsf{X}(\mathrm{A}))\) . For every element \(a\in \mathrm{A}\) , the space \(\mathrm{E}_x\) is stable under \(a\) and if we denote by \(a_{x}\) the endomorphism of \(\mathrm{E}_x\) deduced from a by passing to subspaces, then we have \(a_{x}\circ \theta_{x} = \theta_{x}\circ m_{\mathcal{G}_{\mathrm{A}}(a)}\) .

Let \(\widetilde{\theta}_x\) be the linear map from \(\mathcal{C}(\mathrm{X(A)})\) into \(\mathrm{E}_x\) defined by

\[
\widetilde {\theta} _ {x} (f) = \mathcal {H} _ {\mathrm{A}} (f) (x).
\]

For any function \(f \in \mathcal{C}(\mathsf{X}(\mathsf{A}))\) , we have

\[
\| \widetilde {\theta} _ {x} (f) \| ^ {2} = \langle \mathcal {H} _ {\mathrm{A}} (f) x | \mathcal {H} _ {\mathrm{A}} (f) x \rangle = \langle x | \mathcal {H} _ {\mathrm{A}} (| f | ^ {2}) x \rangle = \int_ {\mathsf {X} (\mathrm{A})} | f | ^ {2} d \mu_ {x}.
\]

Since \(\mathcal{C}(\mathrm{X(A)})\) is dense in \(\mathcal{L}^2 (\mathrm{X(A)},\mu_x)\) , there exists a unique isometric linear map from \(\mathrm{L}^2 (\mathrm{X(A)},\mu_x)\) into \(\mathrm{E}_x\) which extends \(\widetilde{\theta}_x\) ; it is denoted by \(\theta_{x}\) . The map \(\theta_{x}\) is surjective since its image is closed (Lemma 8 of I, p. 107) and contains \(\mathrm{A}\cdot x\) . It is therefore an isometric isomorphism.

Let \(a \in \mathrm{A}\) . Denote by \(g = \mathcal{G}_{\mathrm{A}}(a) \in \mathcal{C}(\mathsf{X}(\mathrm{A}))\) . The space \(\mathrm{E}_x\) is stable under \(a\) . For \(f \in \mathcal{C}(\mathsf{X}(\mathrm{A}))\) , we then have

\[
\widetilde {\theta} _ {x} (m _ {g} (f)) = \mathcal {H} _ {\mathrm{A}} (\mathcal {G} _ {\mathrm{A}} (a) f) x = (a \circ \mathcal {H} _ {\mathrm{A}} (f)) x = a (\widetilde {\theta} _ {x} (f)),
\]

and it follows that \(\theta_x \circ m_g = a_x \circ \theta_x\) .

COROLLARY. — Let E be a complex Hilbert space. Let A be a commutative unital stellar subalgebra of \(\mathcal{L}(\mathrm{E})\) admitting a cyclic vector \(x\) . Let \(\mu_x\) be the spectral measure of \(x\) with respect to A. There exists a unique isometric isomorphism

\[
\theta_ {x} \colon \mathrm{L} ^ {2} (\mathsf {X} (\mathrm{A}), \mu_ {x}) \to \mathrm{E}
\]

such that \(\theta_x(f) = \mathcal{H}_{\mathrm{A}}(f)\) for every function \(f \in \mathcal{C}(\mathsf{X}(\mathrm{A}))\) . For every a in A, we have \(a \circ \theta_x = \theta_x \circ m_{\mathcal{G}_{\mathrm{A}}(a)}\) .

Indeed, we then have \(E_{x}=E\) and \(a_{x}=a\) for all \(a\in A\) .

Let E be a complex Hilbert space. Let \(x \in E\) and let A be a commutative unital stellar subalgebra of \(\mathcal{L}(\mathrm{E})\) . Denote by \(E_{x}\) the closed subspace \(\overline{A \cdot x}\) of E. For every \(a \in A\) , we then have \(a(\mathrm{E}_{x}) \subset \mathrm{E}_{x}\) . Denote by \(A_{x}\) the commutative unital stellar subalgebra of \(\mathcal{L}(\mathrm{E}_{x})\) consisting of the endomorphisms of \(E_{x}\) deduced from the elements of A by passing to subspaces. The vector x is a cyclic vector for \(A_{x} \subset \mathcal{L}(\mathrm{E}_{x})\) .

PROPOSITION 9. --- There exists a subset C of E such that E is the Hilbert sum of the spaces \(E_{x}\) for \(x \in C\) . If E is countably generated, the set C is countable.

Let \(\mathcal{O}\) the set of subsets C of E such that the subspaces \(\mathrm{E}_x\) for \(x\in \mathrm{C}\) are pairwise orthogonal. The set \(\mathcal{O}\) , ordered by inclusion, is of finite character (E, III, p. 34, def. 2) since C belongs to \(\mathcal{O}\) if and only if the sets formed of two elements of C belong to \(\mathcal{O}\) . By E, III, p. 35, th. 1, there exists a maximal element C of \(\mathcal{O}\) .

Let F be the closed subspace of E generated by the subspaces \(E_{x}\) for \(x \in C\) . It suffices to show that \(F^{\circ}\) is reduced to 0 to complete the proof of the proposition. Let y be an element of \(F^{\circ}\) . For every \(x \in C\) and all a and b in A, we have \(\langle a(y) | b(x) \rangle = \langle y | a^{*}b(x) \rangle = 0\) since \(a^{*}b(x)\) belongs to \(E_{x}\) , hence in F. Since the elements \(a(y)\) (resp. \(b(x)\) ) generate a dense subspace of \(E_{y}\) (resp. \(E_{x}\) ), we therefore have \(E_{y} \subset E_{x}^{\circ}\) . Since C is maximal in O, this means that \(E_{y}\) is reduced to 0, hence that y = 0.

Suppose that E is countably generated. Since \(E_{x}\) is nonzero for all \(x \in C\) non-zero, any set C such that E is the Hilbert sum of the spaces \(E_{x}\) for \(x \in C\) is countable.

THEOREM 1. --- Let A be a commutative unital stellar subalgebra of \(\mathcal{L}(\mathrm{E})\) . There exists a locally compact topological space X, a positive measure \(\mu\) on X, an isometric isomorphism \(\theta\) of \(\mathrm{L}^2 (\mathrm{X},\mu)\) onto E and an isometric morphism of stellar algebras \(\pi\) from A into \(\mathcal{C}_b(\mathrm{X})\) such that for every \(a\in \mathrm{A}\) , one has

\[
a \circ \theta = \theta \circ m _ {\pi (a)}.
\]

By prop. 9, there exists a subset C of E such that E is the Hilbert sum of the subspaces \(E_{x}\) for \(x \in C\) . Let X be the locally compact topological space \(\mathsf{X}(\mathsf{A}) \times \mathsf{C}\) , where C is equipped with the discrete topology. There exists a unique measure \(\mu\) on X such that \(\mu|(\mathsf{X}(\mathsf{A}) \times \{x\})\) is identified with the spectral measure \(\mu_{x}\) of x for every \(x \in C\) (INT, III, p. 65, § 2, n°1, prop. 1); this measure is positive (cf. loc. cit.). One then has a Hilbert sum decomposition

\[
\mathrm{L} ^ {2} (\mathrm{X}, \mu) = \bigoplus_ {x \in \mathrm{C}} \mathrm{L} ^ {2} (\mathrm{X} (\mathrm{A}), \mu_ {x})
\]

(prop. 1 of IV, p. 179, c) applied to the sets \(\mathsf{X}(\mathsf{A}) \times \{x\}\) for \(x \in \mathrm{C}\) ). There therefore exists a unique isometry \(\theta: \mathrm{L}^2(\mathrm{X}, \mu) \to \mathrm{E}\) which coincides with \(\theta_x: \mathrm{L}^2(\mathsf{X}(\mathsf{A}), \mu_x) \to \mathrm{E}_x\) on \(\mathrm{L}^2(\mathsf{X}(\mathsf{A}), \mu_x)\) for all \(x \in \mathrm{C}\) .

Let \(p \colon \mathrm{X} \to \mathrm{X}(\mathrm{A})\) be the canonical projection; it is surjective, hence the linear map \(p^* \colon \mathcal{C}_b(\mathrm{X}(\mathrm{A})) \to \mathcal{C}(\mathrm{X})\) defined by \(f \mapsto f \circ p\) is injective. Denote \(\pi = p^* \circ \mathcal{G}_{\mathrm{A}}\) ; this is an injective morphism of stellar algebras from A into \(\mathcal{C}_b(\mathrm{X})\) ; it is therefore isometric (prop. 9 of I, p. 112).

Let \(a \in \mathrm{A}\) . For every \(x \in \mathrm{C}\) , we have \(a_x \circ \theta_x = \theta_x \circ m_{\mathcal{G}_A(a)}\) . The continuous linear maps \(a \circ \theta\) and \(\theta \circ m_{\pi(a)}\) therefore coincide on the subspaces \(\mathrm{L}^2(\mathrm{X}(\mathrm{A}), \mu_x)\) , and are consequently equal.

COROLLARY (Spectral Theorem). --- Let E be a complex Hilbert space. Let \(u \in \mathcal{L}(\mathrm{E})\) be a normal endomorphism of E. There exist a locally compact topological space X, a positive measure \(\mu\) on X, an isometric isomorphism \(\theta\) of \(\mathrm{L}^{2}(\mathrm{X},\mu)\) onto E and a continuous bounded function g on X such that \(u = \theta \circ m_{g} \circ \theta^{-1}\) .

If u admits a cyclic vector x, one may take \(X = \mathrm{Sp}(u)\) and for g the identity function of X.

It suffices to apply Theorem 1 (resp. the corollary of Proposition 8) to the stellar subalgebra A of \(\mathcal{L}(\mathrm{E})\) generated by \(u\) , and to set \(g = \pi(u)\) .

This statement reduces every question concerning a single normal endomorphism of a Hilbert space to a similar question for a multiplication operator, which often simplifies its study (cf. for example Exercise 19 of IV, p. 325).

\subsection{Continuity of the functional calculus}

In this section, we consider the continuity properties of the functional calculus with respect to the two variables.

For a stellar algebra A and a normal element \(a\) of A, and for a function \(f\) continuous on a subset of \(\mathbf{C}\) containing \(\mathrm{Sp}_{\mathrm{A}}(a)\) , we denote \(f(a)\) the element obtained by the continuous functional calculus of \(a\) applied to the restriction of \(f\) to the spectrum of \(a\) .

PROPOSITION 10. --- Let A be a unital stellar algebra. Let U be a relatively compact open set in \(\mathbf{C}\). Denote by \(\Omega_{\mathrm{n}}\) the set of normal elements of A such that \(\mathrm{Sp}_{\mathrm{A}}(a)\) is contained in U. The map \((f,a)\mapsto f(a)\) of \(\mathcal{C}(\overline{\mathrm{U}})\times \Omega_{\mathrm{n}}\) into A is continuous.

Denote by \(q\) the map \((f, a) \mapsto f(a)\) of \(\mathcal{C}(\overline{\mathrm{U}}) \times \Omega_{\mathrm{n}}\) into A. The set of functional calculus maps \(f \mapsto f(a)\) for \(a \in \Omega_{\mathrm{n}}\) is equicontinuous in \(\mathcal{L}(\mathcal{C}(\overline{\mathrm{U}}); \mathrm{A})\) , since each is a continuous linear map of norm \(\leqslant 1\) . To prove the assertion, it therefore suffices to verify that, for every \(f \in \mathcal{C}(\overline{\mathrm{U}})\) , the map from \(\Omega_{\mathrm{n}}\) into A defined by \(a \mapsto f(a)\) is continuous (TG, X, p. 13, cor. 3).

Let \(\mathcal{A}\) be the set of \(f\in\mathcal{C}(\overline{\mathrm{U}})\) such that the map \(a\mapsto f(a)\) of \(\Omega_{\mathrm{n}}\) into A is continuous; it must be proved that \(\mathcal{A}=\mathcal{C}(\overline{\mathrm{U}})\) .

The set \(\mathcal{A}\) is a unital involutive subalgebra of \(\mathcal{C}(\overline{\mathrm{U}})\) . It contains the identity function on \(\overline{\mathrm{U}}\) , hence it separates points. Consequently, it is dense in \(\mathcal{C}(\overline{\mathrm{U}})\) (TG, X, p. 39, prop. 7). Let us prove that it is closed.

Let \((f_{\mathrm{n}})\) be a sequence in \(\mathcal{A}\) converging to \(f \in \mathcal{C}(\overline{\mathrm{U}})\) . Let \(\varepsilon > 0\) and choose \(n \in \mathbf{N}\) such that \(\|f - f_{n}\|_{\infty} \leqslant \varepsilon/4\) . For every \((a_{1}, a_{2}) \in \Omega_{\mathrm{n}}^{2}\) we have

\[
\begin{array}{c} \| f (a _ {1}) - f (a _ {2}) \| \leqslant 2 \| f - f _ {n} \| _ {\infty} + \| f _ {n} (a _ {1}) - f _ {n} (a _ {2}) \| \\ \leqslant \frac {\varepsilon}{2} + \| f _ {n} (a _ {1}) - f _ {n} (a _ {2}) \|. \end{array}
\]

Since \(f_n\) belongs to \(\mathcal{A}\) , there exists a neighborhood V of \(a_1\) in \(\Omega_n\) such that \(\| f_n(a_1) - f_n(a_2) \| \leqslant \varepsilon / 2\) for all \(a_2 \in V\) . We therefore have \(\| f(a_1) - f(a_2) \| \leqslant \varepsilon\) for all \(a_2 \in V\) . The map \(a \mapsto f(a)\) is therefore continuous at \(a_1\) ; the assertion is proved.

COROLLARY 1. --- Let A be a unital stellar algebra and let \(A_{n}\) be the set of normal elements of A. Endow the space \(\mathcal{C}(\mathbf{C})\) with the topology of compact convergence. The map \((f,a)\mapsto f(a)\) of \(\mathcal{C}(\mathbf{C})\times\mathrm{A}_{\mathrm{n}}\) into A is continuous.

Let \((f_{0}, a_{0}) \in \mathcal{C}(\mathbf{C}) \times \mathrm{A}_{\mathrm{n}}\) . Let U be a relatively compact neighborhood of the spectrum of \(a_{0}\) . Let V be the set of normal elements a of A such that \(\mathrm{Sp}_{\mathrm{A}}(a) \subset \mathrm{U}\) ; this is an open neighborhood of \(a_{0}\) in \(A_{n}\) (I, p. 76, prop. 10). For every \(a \in V\) and every function \(f \in \mathcal{C}(\mathbf{C})\) , we have \(f(a) = (f|\overline{\mathrm{U}})(a)\) . Since the map \(f \mapsto \|f|\overline{U}\|_{\infty}\) is a continuous seminorm on \(\mathcal{C}(\mathbf{C})\) endowed with the topology of compact convergence, the continuity of the map \((f, a) \mapsto f(a)\) at \((f_{0}, a_{0})\) follows from prop. 10.

COROLLARY 2. --- Let E be a complex Hilbert space and let \(\mathcal{L}(\mathrm{E})_{\mathrm{n}}\) be the set of normal endomorphisms of E. Endow the space \(\mathcal{C}(\mathbf{C})\) with the topology of compact convergence. The map from \(\mathcal{C}(\mathbf{C}) \times \mathcal{L}(\mathrm{E})_{\mathrm{n}}\) into \(\mathcal{L}(\mathrm{E})\) defined by \((f, u) \mapsto f(u)\) is continuous.

\section{DISTRIBUTIONS AND TEMPERED DISTRIBUTIONS}

In this section, n denotes a natural integer. We denote by \(\mu\) the Lebesgue measure on \(\mathbf{R}^{n}\) , as well as its restriction to every locally compact subset of \(\mathbf{R}^{n}\) .

We denote

\[
x \cdot y = \sum_ {i = 1} ^ {n} x _ {i} y _ {i}
\]

the scalar product on Euclidean space \(\mathbf{R}^{n}\) ; the Euclidean norm is denoted \(x \mapsto \|x\|\) (TG, VI, p. 7). We recall that the group \(\mathbf{R}^{n}\) is in duality with itself with respect to the map \((x, y) \mapsto \exp(2i\pi x \cdot y)\) and that the dual measure of the Lebesgue measure is then identified with the Lebesgue measure (corollary 3 of II, p. 236). We denote by F (resp. \(\overline{F}\) ) the Fourier transform (resp. the inverse Fourier transform) on \(\mathbf{R}^{n}\) (cf. no. 9 of II, p. 237).

For every \(\alpha = (\alpha_{i})_{1\leqslant i\leqslant n}\in \mathbf{N}^{n}\) , we denote by \(X^{\alpha}\) the function from \(\mathbf{R}^n\) into \(\mathbf{R}\) defined by \(x = (x_{i})_{1\leqslant i\leqslant n}\mapsto x^{\alpha} = \prod_{i = 1}^{n}x_{i}^{\alpha_{i}}\) .

Let \(\alpha\) and \(\beta\) be elements of \(\mathbf{N}^n\) . We denote

\[
| \alpha | = \sum_ {i = 1} ^ {n} \alpha_ {i}, \qquad \binom{\alpha}{\beta} = \prod_ {i = 1} ^ {n} \binom{\alpha_ {i}}{\beta_ {i}}.
\]

Lemma 1. --- Let U be an open subset of \(\mathbf{R}^{n}\) .

\begin{enumerate}
\def\labelenumi{\alph{enumi})}
\item
  Let K be a compact subset of U and V an open neighborhood of K in U. There exists a function \(\varphi \in \mathcal{C}^{\infty}(\mathrm{U})\) with compact support contained in V such that \(0 \leqslant \varphi \leqslant 1\) and such that \(\varphi(x) = 1\) for all \(x \in \mathrm{K}\) ;
\item
  For every locally finite open cover of U, there exists a partition of unity consisting of functions of class \(C^{\infty}\) subordinate to it.
\end{enumerate}

This follows from VAR, R1, p. 40, 5.3.6.

PROPOSITION 1. --- Let U be an open subset of \(\mathbf{R}^{n}\) and \(k\in\mathbf{N}\) . Let E, F, and G be topological vector spaces and b: E×F→G a continuous bilinear map. Let f and g be functions of class \(\mathbf{C}^{k}\) from U into E and F, respectively. Then the map h: x↦b(f(x),g(x)) is of class \(\mathbf{C}^{k}\) on U; moreover, for every \(\alpha\in\mathbf{N}^{n}\) such that \(|\alpha|\leqslant k\) and

for all \(x \in \mathrm{U}\) , we have

\[
\partial^{\alpha}h(x) = \sum_{\substack{\beta \in \mathbf{N}^{n}\\ \beta \leqslant \alpha}}\binom {\alpha}{\beta}b\big(\partial^{\beta}f(x),\partial^{\alpha -\beta}g(x)\big).
\]

The map h is the composition of the map \((f,g)\) : \(U \rightarrow E \times F\) , which is of class \(C^{k}\) , and of the map b: \(E \times F \rightarrow G\) which is of class \(C^{\infty}\) . It is therefore of class \(C^{k}\) . The expression for its partial derivatives results from Leibniz's formula (FVR, I, p. 28, prop. 2, cf. A, III, p. 122, corollary).

We recall that a separated locally convex topological vector space is a Montel space if it is barrelled and if every bounded subset is relatively compact (EVT, IV, p. 18, def. 4).

We also recall that if E and F are locally convex spaces, and if E is bornological (EVT, III, p. 12, def. 1), a linear mapping \(u\colon\mathrm{E}\to\mathrm{F}\) is continuous if and only if the image under \(u\) of every bounded subset of E is bounded in F (EVT, III, p. 11, prop. 1, (iiibis), when F is semi-normed, the general case following formally, cf. EVT, II, p. 7, prop. 5, b)).

\subsection{Differentiation under the integral sign}

We fix a locally compact topological space X and a measure \(\nu\) on X. Let E be a Banach space.

PROPOSITION 2. --- Let I be an interval of \(\mathbf{R}\) and f a mapping from X × I into E such that

\begin{enumerate}
\def\labelenumi{(\roman{enumi})}
\item
  For every \(t \in I\) , the map \(x \mapsto f(x, t)\) of X into E is \(\nu\) -integrable;
\item
  For every \(x \in X\) , the map \(t \mapsto f(x, t)\) from I into E admits a derivative, denoted \(t \mapsto f'(x, t)\) ;
\item
  There exists a positive function \(\nu\) -integrable g on X such that \(\|f'(x,t)\| \leqslant g(x)\) for all \((x,t) \in \mathrm{X} \times \mathrm{I}\) .
\end{enumerate}

Then the map F from I into E defined by

\[
\mathrm{F} (t) = \int_ {\mathrm{X}} f (x, t) d \nu (x)
\]

is differentiable in I and for every \(t \in \mathrm{I}\) , we have

\[
\mathrm{F} ^ {\prime} (t) = \int_ {\mathrm{X}} f ^ {\prime} (x, t) d \nu (x).
\]

Let \(t_0 \in \mathrm{I}\) and let \(\mathrm{J}\) be an interval of \(\mathbf{R}\) contained in \(\mathrm{I}\) which is a neighborhood of \(t_0\) in \(\mathrm{I}\) . Let \(h: \mathrm{X} \times \mathrm{J} \to \mathrm{E}\) be the map defined by \((x, t) \mapsto (f(x, t) - f(x, t_0)) / (t - t_0)\) for \((x, t) \in \mathrm{X} \times (\mathrm{J} - \{t_0\})\) and \((x, t_0) \mapsto f'(x, t_0)\) for \(x \in \mathrm{X}\) . Let \(x \in \mathrm{X}\) . We have \(\|h(x, t_0)\| \leqslant g(x)\) and, for all \(t \neq t_0\) , it follows that \(\|h(x, t)\| \leqslant g(x)\) . By FVR, I, p. 23, th. 2. By definition of the derivative, the proposition follows from Corollary 1 of INT, IV, p. 144, § 4, no. 3, applied to the mapping \(h\) .

Let us take up the notations of VAR, R1, 2.4, p. 19 concerning partial derivatives.

COROLLARY 1. --- Let U be an open subset of \(\mathbf{R}^{n}\) . Let \(k \in \mathbf{N}\) and \(f\) be a map from \(\mathrm{X} \times \mathrm{U}\) into E satisfying the following conditions:

\begin{enumerate}
\def\labelenumi{(\roman{enumi})}
\item
  For every \(t \in U\) , the map \(x \mapsto f(x, t)\) of X into E is \(\nu\) -integrable;
\item
  For every \(x \in X\) , the map \(t \mapsto f(x, t)\) from U into E is of class \(C^{k}\) , with partial derivatives denoted \(\partial^{\alpha}f(x, t)\) for all \(\alpha \in \mathbf{N}^{n}\) such that \(|\alpha| \leqslant k\) ;
\item
  There exists a function \(\nu\) -integrable \(g\) on X such that for every \(\alpha \in \mathbf{N}^n\) satisfying \(|\alpha| \leqslant k\) and for every \((x,t) \in \mathrm{X} \times \mathrm{U}\) , we have \(\|\partial^\alpha f(x,t)\| \leqslant g(x)\) .
\end{enumerate}

Then the map F from U to E defined by

\[
\mathrm{F} (t) = \int_ {\mathrm{X}} f (x, t) d \nu (x)
\]

is of class \(C^{k}\) in U and for all \(\alpha\in \mathbf{N}^{n}\) satisfying \(|\alpha|\leqslant k\) and every \(t\in U\) , we have

\[
\partial^ {\alpha} \mathrm{F} (t) = \int_ {\mathrm{X}} \partial^ {\alpha} f (x, t) d \nu (x).
\]

This follows from the proposition by induction on k.

COROLLARY 2. --- Let k be a natural number. Let \(f \in \mathcal{L}^1(\mathbf{R}^n, \mu)\) and \(g \in \mathcal{C}^k(\mathbf{R}^n)\) . Assume that for every \(\alpha \in \mathbf{N}^n\) such that \(|\alpha| \leqslant k\) , the function \(\partial^\alpha g\) is bounded. Then the convolution product \(f * g\) belongs to \(\mathcal{C}^k(\mathbf{R}^n)\) and for every \(\alpha\) such that \(|\alpha| \leqslant k\) , we have \(\partial^\alpha(f * g) = f * \partial^\alpha g\) .

We can apply Corollary 1 to the space \(X = \mathbf{R}^{n}\) with Lebesgue measure and to the mapping h defined by \((x, t) \mapsto f(x)g(t - x)\) from

\(\mathbf{R}^{n} \times \mathbf{R}^{n}\) into \(\mathbf{C}\); indeed, for every \(\alpha \in \mathbf{N}^{n}\) such that \(|\alpha| \leqslant k\) , we have the inequality

\[
| \partial^ {\alpha} h (x, t) | \leqslant \Bigl (\sup _ {| \beta | \leqslant k} \sup _ {y \in \mathbf {R} ^ {n}} | \partial^ {\beta} g (y) | \Bigr) | f (x) |
\]

whose right-hand side is an integrable function on \(R^{n}\) .

\subsection{\texorpdfstring{Integrability criteria in \(\mathbf{R}^{n}\) and \(\mathbf{Z}^{n}\)}{Integrability criteria in \textbackslash mathbf\{R\}\^{}\{n\} and \textbackslash mathbf\{Z\}\^{}\{n\}}}

PROPOSITION 3. --- Let N be a norm on \(\mathbf{R}^{n}\) , \(r\) a real number and \(p \in [1, +\infty[\) .

\begin{enumerate}
\def\labelenumi{\alph{enumi})}
\tightlist
\item
  The function \((1 + \mathrm{N})^{r}\) belongs to \(\mathcal{L}^{p}(\mathbf{R}^{n}, \mu)\) if and only if \(rp < -n\);
\end{enumerate}

a') The restriction to \(\mathbf{Z}^{n}\) of the function \((1+N)^{r}\) belongs to \(\ell^{p}(\mathbf{Z}^{n})\) if and only if \(rp < -n\);

\begin{enumerate}
\def\labelenumi{\alph{enumi})}
\setcounter{enumi}{1}
\tightlist
\item
  Let V be a bounded measurable neighborhood of 0 in \(\mathbf{R}^{n}\) . The function \(\mathbf{N}^{r}\) belongs to \(\mathcal{L}^{p}(\mathbf{R}^{n}-\mathbf{V},\mu)\) if and only if \(rp<-n\) .
\end{enumerate}

By virtue of the equivalence of norms on \(\mathbf{R}^{n}\) (EVT I, p. 14, th. 2 and TG, IX, p. 32, prop. 8), one may assume that the norm N is given by \(\mathrm{N}(x)=\sup|x_{i}|\) for \(x=(x_{i})\in\mathbf{R}^{n}\) . Let B denote the unit ball of \(\mathbf{R}^{n}\) for this norm.

Let V be a bounded measurable neighborhood of 0 in \(\mathbf{R}^n\) . The function \(N^r\) is continuous and bounded on \((\mathrm{B - V})\cup (\mathrm{V - B})\) , which shows that the assertion \(b)\) is valid if and only if it is when \(V = B\) , which we will assume from now on. The function \((1 + N)^r\) being continuous and bounded on B, and satisfying the inequality \(N^r\leqslant (1 + N)^r\leqslant 2^r N^r\) on \(\mathbf{R}^n -\mathbf{B}\) one observes that the assertions \(a)\) and \(b)\) are equivalent.

Let us prove b) when V = B. We may assume that p = 1. The case where \(r > 0\) being elementary, suppose that \(r \leqslant 0\) . For every integer \(j \geqslant 1\) the integrable set

\[
\mathrm{C} _ {j} = \{x \in \mathbf {R} ^ {n} \mid 2 ^ {j - 1} \leqslant \mathrm{N} (x) <   2 ^ {j} \}
\]

has measure \(2^{nj}(2^n - 1)\) . The sets \((\mathrm{C}_j)_{j \geqslant 1}\) form a partition of \(\mathbf{R}^n - \mathbf{B}\) . By INT, V, p. 4, § 1, no. 1, cor., we therefore have

\[
\begin{array}{c} \int_ {\mathbf {R} ^ {n} - \mathrm{B}} ^ {*} \mathrm{N} ^ {r} d \mu = \sum_ {j \geqslant 1} \int_ {\mathrm{C} _ {j}} ^ {*} \mathrm{N} ^ {r} d \mu \\ \leqslant \sum_ {j \geqslant 1} 2 ^ {n + n j} 2 ^ {r (j - 1)} = 2 ^ {n - r} \sum_ {j \geqslant 1} 2 ^ {(n + r) j} \end{array}
\]

which is finite if \(r < -n\) . On the other hand, we have (loc. cit.)

\[
\int_ {\mathbf {R} ^ {n} - \mathrm{B}} ^ {*} \mathrm{N} ^ {r} d \mu \geqslant \sum_ {j \geqslant 1} 2 ^ {r j} 2 ^ {n j} (2 ^ {n} - 1) = (2 ^ {n} - 1) \sum_ {j \geqslant 1} 2 ^ {(r + n) j}
\]

which is infinite if \(r \geqslant -n\) .

One proves \(a'\) ) in a similar manner by considering the sets \(\mathrm{C}_{j} \cap \mathbf{Z}^{n}\) which cover \(\mathbf{Z}^{n} - \{0\}\) and satisfy

\[
\operatorname{Card} \left(\mathrm{C} _ {j} \cap \mathbf {Z} ^ {n}\right) = (2 ^ {j + 1} - 1) ^ {n} - (2 ^ {j} - 1) ^ {n},
\]

which belongs to the interval \([(1-2^{-n})(2^{j+1}-1)^{n},2^{n(j+1)}]\) .

\subsection{Test functions}

In this issue, U denotes an open subset of \(\mathbf{R}^{n}\) . For \(p \in [1, +\infty]\) we will write \(\mathcal{L}^{p}(U)\) and \(L^{p}(U)\) rather than \(\mathcal{L}^{p}(U, \mu)\) and \(L^{p}(U, \mu)\) . We identify the continuous functions belonging to \(\mathcal{L}^{p}(U)\) with their images in \(L^{p}(U)\) .

Equip the space \(\mathcal{C}^{\infty}(\mathrm{U})\) of infinitely differentiable functions in U with complex values, with the topology defined by the family of seminorms \(p_{\alpha,K}\) defined for \(\alpha \in \mathbf{N}^{n}\) and \(K \subset U\) by

\[
p _ {\alpha , \mathrm{K}} (\varphi) = \sup _ {x \in \mathrm{K}} | \partial^ {\alpha} \varphi (x) |.
\]

This space is complete (cf. EVT, III, p. 9, example b)).

For every compact subset K of U, we denote by \(\mathcal{C}_{\mathrm{K}}^{\infty}(\mathrm{U})\) the subspace of \(\mathcal{C}^{\infty}(\mathrm{U})\) formed by functions with support in K. The space \(\mathcal{C}_{\mathrm{K}}^{\infty}(\mathrm{U})\) is equipped with the topology induced by that of \(\mathcal{C}^{\infty}(\mathrm{U})\) . It is a Fréchet space, a countable family of seminorms defining its topology being the family \((p_{\alpha,\mathrm{K}})_{\alpha\in\mathbf{N}^{n}}\) . In particular, it is a bornological space (EVT, III, p. 12, prop. 2).

We denote by \(\mathcal{D}(\mathrm{U})\) the vector space \(\mathcal{K}(\mathrm{U})\cap \mathcal{C}^{\infty}(\mathrm{U})\) of functions of class \(\mathrm{C}^{\infty}\) with compact support in U. We say that \(\mathcal{D}(\mathrm{U})\) is the space of test functions in U. The space \(\mathcal{D}(\mathrm{U})\) is the inductive limit of the spaces \(\mathcal{C}_{\mathrm{K}}^{\infty}(\mathrm{U})\) and it is equipped with the corresponding locally convex inductive limit topology (EVT, II, p. 31).

This space is denoted \(\mathcal{C}_{\circ}^{\infty}(U)\) in EVT, III, p. 9.

Let \((\mathrm{K}_{m})\) be an increasing sequence of compact subsets of U whose interiors form a covering of U. The space \(\mathcal{D}(\mathrm{U})\) is then the strict inductive limit of the spaces \(\mathcal{C}_{\mathrm{K}_{m}}^{\infty}(\mathrm{U})\) (EVT, III, p. 9). It is therefore a complete space (EVT, II, p. 35, prop. 9). Every bounded subset of \(\mathcal{D}(\mathrm{U})\) is contained in one of the subspaces \(\mathcal{C}_{\mathrm{K}_{m}}^{\infty}(\mathrm{U})\) (EVT, III, p. 5, prop. 6). This means that \(\mathrm{B}\subset\mathcal{D}(\mathrm{U})\) is bounded if and only if there exists a compact subset K of U and a family \((\mathrm{M}_{\alpha})_{\alpha\in\mathbf{N}^{n}}\) in \(R_{+}\) such that B is contained in the set of functions \(\varphi\in\mathcal{C}_{\mathrm{K}}^{\infty}(\mathrm{U})\) satisfying

\[
p _ {\alpha , \mathrm{K}} (\varphi) \leqslant \mathrm{M} _ {\alpha}
\]

for all \(\alpha\in\mathbf{N}^{n}\) .

The space \(\mathcal{D}(\mathrm{U})\) is a bornological space (EVT, III, p. 12, example 3) and a Montel space (EVT, IV, p. 18, example 4).

Let V be an open subset of \(\mathbf{R}^{n}\) contained in U. If \(K \subset V\) is compact, then the restriction of functions to V defines a continuous and surjective linear map from \(\mathcal{C}_{\mathrm{K}}^{\infty}(\mathrm{U})\) onto \(\mathcal{C}_{\mathrm{K}}^{\infty}(\mathrm{V})\) . The extension by zero of a function defined on V induces an injective continuous linear map, called canonical, from \(\mathcal{D}(\mathrm{V})\) into \(\mathcal{D}(\mathrm{U})\) .

Remark. --- One defines similarly the space \(\mathcal{D}_{\mathbf{R}}(\mathrm{U})\) of test functions in U with real values. The linear map such that \(z \otimes \varphi \mapsto z\varphi\) for all \(z \in \mathbf{C}\) and every \(\varphi \in \mathcal{D}_{\mathbf{R}}(\mathrm{U})\) is an isomorphism of \(\mathbf{C} \otimes \mathcal{D}_{\mathbf{R}}(\mathrm{U})\) into \(\mathcal{D}(\mathrm{U})\) .

Lemma 2. --- Let \(\mathcal{U}\) be an open covering of U. There exists a locally finite open covering of U finer than \(\mathcal{U}\) and consisting of relatively compact sets.

Since U is locally compact, there exists a covering \(\mathcal{V}\) of U finer than \(\mathcal{U}\) consisting of open sets relatively compact in U. Since U is paracompact (TG, IX, p. 51, th. 4), there exists a locally finite open covering \(\mathcal{W}\) finer than \(\mathcal{V}\). Then \(\mathcal{W}\) is finer than \(\mathcal{U}\) and consists of relatively compact open sets.

Lemma 3. --- Let \(f \in \mathcal{K}(\mathbf{R}^n)\) and let V be an open neighborhood of the support K of f. There exists an integer \(m_0 \geqslant 1\) such that, for every \(x \in K\) , the set V contains the ball of radius \(m_0^{-1}\) centered at x. For every integer \(m > m_0\) , let \(V_m\) be the closed ball of radius \(m^{-1}\) centered at 0 in \(\mathbf{R}^{n}\) and let \(\varphi_{m}\) be a positive test function with support in \(V_{m}\) and integral 1. Set \(f_{m} = \varphi_{m} * f\) .

\begin{enumerate}
\def\labelenumi{\alph{enumi})}
\item
  One has \(f_{m}\in \mathcal{D}(\mathbf{R}^{n})\) and the support of \(f_{m}\) is contained in V;
\item
  Let \(p \in [1, +\infty]\) . The sequence \((f_m)_{m > m_0}\) converges to \(f\) in \(L^p(\mathbf{R}^n)\) .
\end{enumerate}

According to TG, IX, p. 14, remark, we have \(d(\mathrm{K}, \mathbf{R}^{n}-\mathrm{V}) > 0\) . Let \(m_{0} \geqslant 1\) be an integer such that \(m_{0}^{-1} < d(\mathrm{K}, \mathbf{R}^{n}-\mathrm{V})\) ; it satisfies the required condition. For all \(m > m_{0}\) , the function \(f_{m}\) is continuous (INT, VIII, p. 166, § 4, n° 5, prop. 11) and has support contained in \(\mathrm{K} + \mathrm{V}_{m}\) (INT, VIII, p. 126, § 1, n° 4, prop. 5), hence in V. By Corollary 2 of IV, p. 198, we have \(f_{m} \in \mathcal{D}(\mathbf{R}^{n})\) . If \(p \neq +\infty\) , then the sequence ( \(f_{m}\) ) converges to \(f\) in \(\mathrm{L}^{p}(\mathbf{R}^{n})\) (INT, VIII, p. 172, § 4, no. 7, prop. 20). \(^{(1)}\)

Suppose that \(p = +\infty\) . Let \(\varepsilon > 0\) . The function \(f\) is uniformly continuous on \(\mathbf{R}^n\) ; therefore there exists an integer \(m_1 > m_0\) such that, for every \(m \geqslant m_1\) and for all \(x \in \mathbf{R}^n\) and \(y \in \mathrm{V}_m\) , one has \(|f(x) - f(x - y)| < \varepsilon\) .

For every \(m \geqslant m_{1}\) , the function \(\varphi_{m}\) vanishes outside \(\mathrm{V}_{m}\) , is positive, and has integral 1. We therefore deduce the inequality

\[
| f (x) - f _ {m} (x) | = \left| \int_ {\mathbf {R} ^ {n}} (f (x) - f (x - y)) \varphi_ {m} (y) d \mu (y) \right| \leqslant \varepsilon
\]

for all \(x \in \mathbf{R}^{n}\) , whence the result.

Proposition 4. --- a) The canonical injection of \(\mathcal{D}(\mathrm{U})\) into \(\mathcal{K}(\mathrm{U})\) is continuous and \(\mathcal{D}(\mathrm{U})\) is dense in \(\mathcal{K}(\mathrm{U})\) ;

\begin{enumerate}
\def\labelenumi{\alph{enumi})}
\setcounter{enumi}{1}
\tightlist
\item
  For every \(p \in [1, +\infty[\) , the canonical injection of \(\mathcal{D}(\mathrm{U})\) into \(\mathcal{L}^p(\mathrm{U})\) is continuous and \(\mathcal{D}(\mathrm{U})\) is dense in \(\mathrm{L}^p(\mathrm{U})\) .
\end{enumerate}

Every continuous seminorm on \(\mathcal{K}(\mathrm{U})\) is continuous on \(\mathcal{D}(\mathrm{U})\) , hence the canonical injection of \(\mathcal{D}(\mathrm{U})\) into \(\mathcal{K}(\mathrm{U})\) is continuous.

Let \(f \in \mathcal{K}(\mathrm{U})\) . There exists a sequence \((f_{m})_{m \in \mathbf{N}}\) in \(\mathcal{D}(\mathrm{U})\) which converges uniformly to f on U and such that the supports of the \(f_{m}\) are contained in a fixed relatively compact neighborhood of the support of f (Lemma 3). The sequence \((f_{m})\) therefore converges to f in \(\mathcal{K}(\mathrm{U})\) . This concludes the proof of a).

Let \(p \in [1, +\infty[\) . Let \(f \in \mathcal{D}(U)\) and let K be the support of f. We then have the inequality \(\mathrm{N}_{p}(f) \leqslant \mu(\mathrm{K})^{1/p} \sup_{x \in \mathrm{K}} |f(x)|\) , hence the canonical injection of \(\mathcal{D}(U)\) into \(\mathcal{L}^{p}(U)\) is continuous. The last part of assertion b) then follows from a).

The product of two test functions being again a test function, the Leibniz formula (FVR, I, p. 28, prop. 2) shows that the space \(\mathcal{D}(\mathrm{U})\) is a topological algebra.

More generally, let \(f \in \mathcal{C}^{\infty}(\mathrm{U})\) . The linear map \(\varphi \mapsto f\varphi\) from \(\mathcal{D}(\mathrm{U})\) into \(\mathcal{D}(\mathrm{U})\) is then continuous. Indeed, for every compact subset \(\mathbf{K}\) of \(\mathrm{U}\) and every \(\alpha \in \mathbf{N}^{n}\) , we have

\[
p _ {\alpha , \mathrm{K}} (f \varphi) \leqslant \sum_ {0 \leqslant \beta \leqslant \alpha} \binom {\alpha} {\beta} p _ {\beta , \mathrm{K}} (f) p _ {\alpha - \beta , \mathrm{K}} (\varphi)
\]

(cf. loc. cit.).

\subsection{Distributions}

We keep the notations of the preceding number; U therefore denotes an open subset of \(R^{n}\) .

DEFINITION 1. --- The dual space of \(\mathcal{D}(\mathrm{U})\) , endowed with the topology of bounded convergence, is called the space of distributions on U. It is denoted by \(\mathcal{D}'(\mathrm{U})\) .

A distribution on U is therefore a continuous linear form on \(\mathcal{D}(\mathrm{U})\) .

If \(f\) is a linear form on \(\mathcal{D}(\mathrm{U})\) (and in particular if \(f\) is a distribution) and if \(\varphi\) belongs to \(\mathcal{D}(\mathrm{U})\) , we denote by \(\langle f, \varphi \rangle\) the evaluation of \(f\) at \(\varphi\) .

Since \(\mathcal{D}(\mathrm{U})\) is bornological, the space \(\mathcal{D}'(\mathrm{U})\) is complete (TVS, III, p. 24, cor. 1). Since \(\mathcal{D}(\mathrm{U})\) is a Montel space, the same holds for \(\mathcal{D}'(\mathrm{U})\) (EVT, IV, p. 19, prop. 9). In particular, the space \(\mathcal{D}'(\mathrm{U})\) is reflexive (EVT, IV, p. 16, th. 2).

Lemma 4. --- Let f be a linear map from \(\mathcal{D}(\mathrm{U})\) into \(\mathbf{C}\) . Then \(f\) is a distribution if and only if, for every compact subset \(\mathrm{K}\) of \(\mathrm{U}\) and for every family \((\mathrm{M}_{\alpha})_{\alpha \in \mathbf{N}^{n}}\) in \(\mathbf{R}_{+}\) , the linear form \(f\) is bounded on the set of functions \(\varphi \in \mathcal{C}_{\mathrm{K}}^{\infty}(\mathrm{U})\) such that, for every \(\alpha \in \mathbf{N}^{n}\) , we have

\[
\sup _ {x \in \mathrm{K}} | \partial^ {\alpha} \varphi (x) | \leqslant \mathrm{M} _ {\alpha}.
\]

Indeed, the space \(\mathcal{D}(\mathrm{U})\) is bornological and every bounded subset of \(\mathcal{D}(\mathrm{U})\) is contained in one of the bounded sets described in the statement.

Lemma 5. --- Let F be a filter on \(\mathcal{D}'(\mathrm{U})\) having a countable base or containing a simply bounded set. Then F converges to a distribution if and only if \(\langle f, \varphi \rangle\) converges according to F for every test function \(\varphi\) on U.

In particular, a sequence \((f_{m})_{m\in\mathbf{N}}\) of distributions converges to a distribution f if and only if, for every \(\varphi\in\mathcal{D}(\mathrm{U})\) the sequence \((\langle f_{m},\varphi\rangle)_{m\in\mathbf{N}}\) converges to \(\langle f,\varphi\rangle\) .

Since \(\mathcal{D}'(\mathrm{U})\) is a Montel space, hence barrelled; this follows from the Banach–Steinhaus theorem (EVT, III, p. 26, cor. 3 of th. 1).

Let V be an open set contained in U. The transpose of the canonical linear map from \(\mathcal{D}(\mathrm{V})\) into \(\mathcal{D}(\mathrm{U})\) is a continuous linear map from \(\mathcal{D}'(\mathrm{U})\) into \(\mathcal{D}'(\mathrm{V})\) , called the restriction of the distributions from U to V and denoted \(r_{\mathrm{VU}}\) , or sometimes \(r_{\mathrm{V,U}}\) .

One has \(r_{\mathrm{VV}} = 1_{\mathcal{D}'(\mathrm{V})}\) . If \(W \subset V \subset U\) are open subsets of U, then we have \(r_{\mathrm{WV}} \circ r_{\mathrm{VU}} = r_{\mathrm{WU}}\) . In other words, if \(\mathcal{T}\) denotes the topology on U, the projective system \(\underline{\mathcal{D}'}(\mathrm{U}) = ((\mathcal{D}'(\mathrm{V}))_{\mathrm{V} \in \mathcal{T}}, (r_{\mathrm{WV}}))\) is a presheaf on U with values in the species of structure of locally convex topological vector spaces (TA, I, p. 42, def. 1 and p. 66, §10).

PROPOSITION 5. --- The presheaf \(\underline{\mathcal{D}}^{\prime}(\mathrm{U})\) is a sheaf.

Recall (TA, I, p. 43, def. 2) that this means that for every open subset V of U and every open covering \((\mathrm{V}_{i})_{i\in\mathrm{I}}\) of V, the following conditions are satisfied:

\begin{enumerate}
\def\labelenumi{(\roman{enumi})}
\item
  The map \((r_{\mathrm{V}_i\mathrm{V}})_{i\in \mathrm{I}}\colon \mathcal{D}'(\mathrm{V})\to \prod_{i\in \mathrm{I}}\mathcal{D}'(\mathrm{V}_i)\) is injective;
\item
  For every family \((f_i)\in \prod_{i\in \mathrm{I}}\mathcal{D}'(\mathrm{V}_i)\) such that
\end{enumerate}

\[
r _ {(\mathrm{V} _ {i} \cap \mathrm{V} _ {i ^ {\prime}}) \mathrm{V} _ {i}} (f _ {i}) = r _ {(\mathrm{V} _ {i} \cap \mathrm{V} _ {i ^ {\prime}}) \mathrm{V} _ {i ^ {\prime}}} (f _ {i ^ {\prime}})
\]

for every pair \((i, i') \in \mathrm{I} \times \mathrm{I}\) , there exists a distribution \(f \in \mathcal{D}'(\mathrm{V})\) such that for every \(i \in I\) , one has \(r_{\mathrm{V}_{i}\mathrm{V}}(f) = f_{i}\) .

Let V be an open subset of U and \(\mathcal{V} = (\mathrm{V}_{i})_{i \in \mathrm{I}}\) an open covering of V. Let \((\mathrm{W}_{j})_{j \in \mathrm{J}}\) be an open covering of V that is locally finite, finer than \(\mathcal{V}\), and consisting of relatively compact parts (lemma 2 of IV, p. 201). Fix a partition of unity \((\varphi_{j})_{j \in \mathrm{J}}\) subordinate to the covering \((\mathrm{W}_{j})_{j \in \mathrm{J}}\) such that the support of \(\varphi_{j}\) is contained in \(W_{j}\) for all \(j \in J\) (lemma 1 of IV, p. 196).

Let \(\varphi\in\mathcal{D}(V)\) . Since the set of \(j\in J\) such that \(W_{j}\) meets the support of \(\varphi\) is finite (TG, I, §9, no. 1, p. 59), we have

\[
\varphi = \sum_ {j \in \mathrm{J}} \varphi \varphi_ {j}
\]

where the sum has only finitely many nonzero terms.

Let us prove (i). Suppose that \(f \in \mathcal{D}'(\mathrm{V})\) satisfies \(r_{\mathrm{V}_i\mathrm{V}}(f) = 0\) for all \(i \in \mathrm{I}\) . This means that \(\langle f, \varphi \rangle = 0\) for every test function \(\varphi\) whose support is contained in one of the open sets \(\mathrm{V}_i\) . This is a fortiori true if the support of \(\varphi\) is contained in one of the open sets \(\mathrm{W}_j\) . But then, for every \(\varphi \in \mathcal{D}(\mathrm{V})\) , we have

\[
\langle f, \varphi \rangle = \langle f, \sum_ {j \in \mathrm{J}} \varphi \varphi_ {j} \rangle = \sum_ {j \in \mathrm{J}} \langle f, \varphi \varphi_ {j} \rangle = 0,
\]

hence f = 0.

Let us prove (ii). Let \((f_{i})_{i\in\mathrm{I}}\) be a family such that \(f_{i}\in\mathcal{D}^{\prime}(\mathrm{V}_{i})\) for all \(i\in I\) and \(r_{\mathrm{V}_{i}\cap\mathrm{V}_{i^{\prime}},\mathrm{V}_{i}}(f_{i})=r_{\mathrm{V}_{i}\cap\mathrm{V}_{i^{\prime}},\mathrm{V}_{i^{\prime}}}(f_{i^{\prime}})\) for all i and \(i^{\prime}\) in I. Let \(\iota:J\to I\) be a map such that \(\mathrm{W}_{j}\subset\mathrm{V}_{\iota(j)}\) for all \(j\in J\) . For every \(j\in J\) set \(\widetilde{f}_{j}=r_{\mathrm{W}_{j},\mathrm{V}_{\iota(j)}}(f_{\iota(j)})\) . We then have

\[
\begin{array}{r l} & r _ {\mathrm{W} _ {j} \cap \mathrm{W} _ {j ^ {\prime}}, \mathrm{W} _ {j}} (\widetilde {f} _ {j}) = r _ {\mathrm{W} _ {j} \cap \mathrm{W} _ {j ^ {\prime}}, \mathrm{W} _ {j}} (r _ {\mathrm{W} _ {j}, \mathrm{V} _ {\iota (j)}} (f _ {\iota (j)})) \\ & \qquad = r _ {\mathrm{W} _ {j} \cap \mathrm{W} _ {j ^ {\prime}}, \mathrm{V} _ {\iota (j)}} (f _ {\iota (j)}) \\ & \qquad = r _ {\mathrm{W} _ {j} \cap \mathrm{W} _ {j ^ {\prime}}, \mathrm{V} _ {\iota (j)} \cap \mathrm{V} _ {\iota (j ^ {\prime})}} \circ r _ {\mathrm{V} _ {\iota (j)} \cap \mathrm{V} _ {\iota (j ^ {\prime})}, \mathrm{V} _ {\iota (j)}} (f _ {\iota (j)}). \end{array}
\]

By exchanging the roles of j and j' and noting that

\[
r _ {\mathrm{V} _ {\iota (j)} \cap \mathrm{V} _ {\iota (j ^ {\prime})}, \mathrm{V} _ {\iota (j)}} (f _ {\iota (j)}) = r _ {\mathrm{V} _ {\iota (j)} \cap \mathrm{V} _ {\iota (j ^ {\prime})}, \mathrm{V} _ {\iota (j ^ {\prime})}} (f _ {\iota (j ^ {\prime})}),
\]

by hypothesis, we deduce that

\[
r _ {\mathrm{W} _ {j} \cap \mathrm{W} _ {j ^ {\prime}}, \mathrm{W} _ {j}} (\widetilde {f} _ {j}) = r _ {\mathrm{W} _ {j} \cap \mathrm{W} _ {j ^ {\prime}}, \mathrm{W} _ {j}} (\widetilde {f} _ {j ^ {\prime}})\tag{1}
\]

for all \((j, j') \in \mathrm{J}^{2}\) .

For \(\varphi \in \mathcal{D}(\mathrm{V})\) , set

\[
\lambda (\varphi) = \sum_ {j \in \mathrm{J}} \langle \widetilde {f} _ {j}, \varphi \varphi_ {j} | \mathrm{W} _ {j} \rangle ,
\]

where the sum is finite since only finitely many terms can be nonzero.

The map \(\lambda\) is a linear form on \(\mathcal{D}(\mathrm{V})\) . Let us prove that it is a distribution. Let B be a bounded subset of \(\mathcal{D}(\mathrm{V})\) and let K be a compact subset of V such that \(B \subset \mathcal{C}_{K}^{\infty}(V)\) . By TG, I, §9, no. 1, p. 59, there exists a finite subset J' of J such that

\[
\lambda (\varphi) = \sum_ {j \in J ^ {\prime}} \left\langle \widetilde {f} _ {j}, \varphi \varphi_ {j} \mid W _ {j} \right\rangle
\]

for all \(\varphi\in\mathrm{B}\) . Since \(\widetilde{f}_{j}\) is a distribution for every \(j\in\mathrm{J}^{\prime}\) and that \(\mathcal{D}(\mathrm{V})\) is a topological algebra, it follows that \(\lambda\) is bounded on B, whence the result (Lemma 4).

Let \(j \in J\) and let \(\varphi \in \mathcal{D}(V)\) have support contained in \(W_{j}\) . We then have

\[
\langle \lambda , \varphi \rangle = \sum_ {j ^ {\prime} \in \mathrm{J}} \left\langle \widetilde {f} _ {j ^ {\prime}}, \varphi \varphi_ {j ^ {\prime}} \right| \mathrm{W} _ {j} \rangle = \sum_ {j ^ {\prime} \in \mathrm{J}} \left\langle \widetilde {f} _ {j}, \varphi \varphi_ {j ^ {\prime}} \right| \mathrm{W} _ {j} \rangle
\]

according to (1), since \(\varphi\varphi_{j'}\) has support contained in \(W_j \cap W_{j'}\) . Consequently

\[
\langle \lambda , \varphi \rangle = \langle \widetilde {f} _ {j}, \sum_ {j ^ {\prime} \in \mathrm{J}} \varphi \varphi_ {j ^ {\prime}} | \mathrm{W} _ {j} \rangle = \langle \widetilde {f} _ {j}, \varphi | \mathrm{W} _ {j} \rangle ,
\]

whence \(r_{\mathrm{W}_j\mathrm{V}}(\lambda) = \widetilde{f}_j\) for all \(j\in \mathbf{J}\) .

Let \(i \in I\) . Let us finally prove that the restriction of \(\lambda\) to \(V_i\) coincides with \(f_i\) . By condition (i), applied to the covering of \(V_i\) by the open sets \(W_j\) , it suffices to verify that for every \(j \in J\) the restriction of \(\lambda\) to \(V_i \cap W_j\) coincides with that of \(f_i\) . By the preceding, it remains to verify that the restriction of \(f_i\) to \(V_i \cap W_j\) is that of \(\widetilde{f}_j\) . Now, we have

\[
\begin{array}{c} r _ {\mathrm{V} _ {i} \cap \mathrm{W} _ {j}, \mathrm{V} _ {i}} (f _ {i}) = r _ {\mathrm{V} _ {i} \cap \mathrm{W} _ {j}, \mathrm{V} _ {i} \cap \mathrm{V} _ {\iota (j)}} (r _ {\mathrm{V} _ {i} \cap \mathrm{V} _ {\iota (j)}, \mathrm{V} _ {i}} (f _ {i})) \\ = r _ {\mathrm{V} _ {i} \cap \mathrm{W} _ {j}, \mathrm{V} _ {i} \cap \mathrm{V} _ {\iota (j)}} (r _ {\mathrm{V} _ {i} \cap \mathrm{V} _ {\iota (j)}, \mathrm{V} _ {\iota (j)}} (f _ {\iota (j)})) = r _ {\mathrm{V} _ {i} \cap \mathrm{W} _ {j}, \mathrm{V} _ {\iota (j)}} (f _ {\iota (j)}), \end{array}
\]

where we used the hypothesis concerning the family \((f_{i})\) . But then

\[
r _ {\mathrm{V} _ {i} \cap \mathrm{W} _ {j}, \mathrm{V} _ {\iota (j)}} (f _ {\iota (j)}) = r _ {\mathrm{V} _ {i} \cap \mathrm{W} _ {j}, \mathrm{W} _ {j}} (r _ {\mathrm{W} _ {j}, \mathrm{V} _ {\iota (j)}} (f _ {\iota (j)})) = r _ {\mathrm{V} _ {i} \cap \mathrm{W} _ {j}, \mathrm{W} _ {j}} (\widetilde {f} _ {j}),
\]

which allows us to conclude.

\subsection{Interpretation of functions as distributions}

PROPOSITION 6. --- Let ν be a measure on U. The restriction of ν to D(U) is a distribution, which is zero if and only if the measure ν is zero.

Let K be a compact subset of U. For every function \(\varphi\in\mathcal{C}_{\mathrm{K}}^{\infty}(\mathrm{U})\) , we have \(|\langle\nu,\varphi\rangle|\leqslant p_{0,\mathrm{K}}(\varphi)|\nu|(\mathrm{K})\) , hence the restriction of \(\nu\) to \(\mathcal{D}(\mathrm{U})\) is continuous. Since \(\mathcal{D}(\mathrm{U})\) is dense in \(\mathcal{K}(\mathrm{U})\) , by proposition 4(a) of IV, p. 202, the restriction of \(\nu\) to \(\mathcal{D}(\mathrm{U})\) is zero if and only if \(\nu\) is zero.

We will identify the space \(\mathcal{M}(\mathrm{U};\mathbf{C})\) of complex measures on U with a subspace of \(\mathcal{D}'(\mathrm{U})\) . We will also identify the space \(\mathrm{L}_{\mathrm{loc}}^{1}(\mathrm{U})\) with a subspace of \(\mathcal{D}'(\mathrm{U})\) by the map induced by passing to quotients of the map that assigns to \(f\in \mathcal{L}_{\mathrm{loc}}^{1}(\mathrm{U})\) the measure \(f\cdot \mu\) (INT, V, p. 44, § 5, no. 2, def. 2) . In other words, for \(f\in \mathrm{L}_{\mathrm{loc}}^{1}(\mathrm{U})\) and \(\varphi \in \mathcal{D}(\mathrm{U})\) , we have

\[
\langle f, \varphi \rangle = \int_ {\mathrm{U}} f \varphi d \mu .
\]

In particular, for \(p \in [1, +\infty]\) this allows one to identify the space \(\mathrm{L}^p(\mathrm{U})\) with a subspace of \(\mathcal{D}'(\mathrm{U})\) (cf. INT, V, p. 43, § 5, no. 1).

PROPOSITION 7. --- Let \(p \in [1, +\infty]\) . The injection of \(\mathrm{L}^p(\mathrm{U})\) into \(\mathcal{D}'(\mathrm{U})\) is continuous.

Let \(f \in \mathrm{L}^{p}(\mathrm{U})\) . Let K be a compact subset of U, and denote by \(\varphi_{K}\) its characteristic function. For every test function \(\varphi\) with support contained in K, Hölder's inequality implies

\[
| \langle f, \varphi \rangle | = \left| \int_ {\mathrm{U}} f \varphi d \mu \right| \leqslant \mathrm{N} _ {p} (f) \mathrm{N} _ {q} (\varphi) \leqslant \mathrm{N} _ {p} (f) \mathrm{N} _ {q} (\varphi_ {\mathrm{K}}) p _ {0, \mathrm{K}} (\varphi)
\]

where q is the conjugate exponent of p.

In particular, one can identify \(\mathcal{D}(\mathrm{U})\) with a subspace of \(\mathcal{D}'(\mathrm{U})\) .

PROPOSITION 8. --- The space \(\mathcal{D}(\mathrm{U})\) is dense in \(\mathcal{D}'(\mathrm{U})\) .

Let \(\lambda\) be a linear form on \(\mathcal{D}'(\mathrm{U})\) vanishing on \(\mathcal{D}(\mathrm{U})\) . Since \(\mathcal{D}(\mathrm{U})\) is reflexive, there exists a test function \(\varphi \in \mathcal{D}(\mathrm{U})\) such that \(\lambda(f) = \langle f, \varphi \rangle\) for all \(f \in \mathcal{D}'(\mathrm{U})\) . We obtain

\[
0 = \lambda (\overline {{\varphi}}) = \langle \overline {{\varphi}}, \varphi \rangle = \int_ {\mathrm{U}} | \varphi | ^ {2} d \mu ,
\]

whence \(\varphi = 0\) . The proposition then follows from the Hahn--Banach theorem (EVT, II, p. 49, cor. 3 (ii)).

For \(h \in \mathcal{C}^{\infty}(\mathrm{U})\) , the transpose of the continuous linear map \(\varphi \mapsto h\varphi\) of \(\mathcal{D}(\mathrm{U})\) into itself is a continuous linear map from \(\mathcal{D}'(\mathrm{U})\) into itself, which is denoted \(f \mapsto hf\) . This definition is justified because if \(f\) is the distribution associated with a measure \(\nu\) on \(\mathrm{U}\) , then \(hf\) is associated with the measure \(h \cdot \nu\) . Indeed, for every test function \(\varphi \in \mathcal{D}(\mathrm{U})\) , one computes

\[
\langle h f, \varphi \rangle = \langle f, h \varphi \rangle = \int_ {\mathrm{U}} h \varphi d \nu = \int_ {\mathrm{U}} \varphi d (h \cdot \nu).
\]

\subsection{Differentiation of distributions}

Let \(\alpha \in \mathbf{N}^n\) . The linear map \(\varphi \mapsto \partial^\alpha \varphi\) is continuous from \(\mathcal{D}(\mathrm{U})\) into \(\mathcal{D}(\mathrm{U})\) . Its transpose \(^t\partial^\alpha\) is a continuous linear map from \(\mathcal{D}'(\mathrm{U})\) into \(\mathcal{D}'(\mathrm{U})\) (EVT, IV, p. 6, cor., \(b\) )).

We denote \(\partial^{\alpha}\) the continuous linear map \((-1)^{|\alpha|}{}^{t}\partial^{\alpha}\) from \(\mathcal{D}'(\mathrm{U})\) into itself.

DEFINITION 2. --- If \(f \in \mathcal{D}'(\mathrm{U})\) is a distribution, one says that \(\partial^{\alpha} f\) is the iterated partial derivative of order \(\alpha\) of \(f\) .

Thus, by definition

\[
\langle \partial^ {\alpha} f, \varphi \rangle = (- 1) ^ {| \alpha |} \langle f, \partial^ {\alpha} \varphi \rangle
\]

for every function \(\varphi\in\mathcal{D}(\mathrm{U})\) . We have \(\partial^{\alpha+\beta}=\partial^{\alpha}\circ\partial^{\beta}\) for all \(\alpha\) and \(\beta\) in \(\mathbf{N}^{n}\) .

If \(n=1\) , one also denotes \(f'\) the derivative of a distribution \(f\in\mathcal{D}'(U)\) . The definition is justified by the following lemma.

Lemma 6. --- Let k be a natural integer. Let \(f \in \mathcal{C}^k(\mathrm{U})\) and let \(\lambda\) be the distribution associated with \(f\) . For every \(\beta \in \mathbf{N}^n\) such that \(|\beta| \leqslant k\) , the distribution \(\partial^\beta\lambda\) is the distribution associated with the function \(\partial^\beta f\) .

By induction on \(k\) , it suffices to prove this property when \(\beta\) satisfies \(|\beta| = 1\) , and one may even assume that \(\beta = (0, \ldots, 0, 1)\) . Since distributions form a sheaf (prop. 5 of IV, p. 204), it suffices to verify the assertion when there exists an open set \(V \subset \mathbf{R}^{n-1}\) and an open interval \(I \subset \mathbf{R}\) such that \(U = V \times I\) .

For every test function \(\varphi\in\mathcal{D}(\mathrm{U})\) , one has by definition

\[
\langle \partial^ {\beta} \lambda , \varphi \rangle = - \int_ {\mathrm{U}} f (x) \partial^ {\beta} \varphi (x) d x = - \int_ {\mathrm{V}} \Bigl (\int_ {\mathrm{I}} f (y, t) \partial^ {\beta} \varphi (y, t) d t \Bigr) d y
\]

by the Lebesgue--Fubini theorem (INT, V, p. 96, § 8, n° 4, th. 1). By integration by parts (FVR, II, p. 10), one has

\[
- \int_ {\mathrm{I}} f (y, t) \partial^ {\beta} \varphi (y, t) d t = \int_ {\mathrm{I}} \partial^ {\beta} f (y, t) \varphi (y, t) d t,
\]

since \(t \mapsto \varphi(y, t)\) is compactly supported in I. One therefore obtains

\[
\langle \partial^ {\beta} \lambda , \varphi \rangle = \int_ {V} \left(\int_ {I} \partial^ {\beta} f (y, t) \varphi (y, t) d t\right) d y = \langle \partial^ {\beta} f, \varphi \rangle .
\]

PROPOSITION 9 (Leibniz formula). --- Let f be a distribution on U and g an infinitely differentiable function on U. Let \(\alpha \in \mathbf{N}^n\) . One has the relation

\[
\partial^ {\alpha} (f g) = \sum_ {\beta \leqslant \alpha} \binom {\alpha} {\beta} \partial^ {\beta} f   \partial^ {\alpha - \beta} g.
\]

Proceeding by induction on \(|\alpha|\) as in the proof of FVR, I, p. 28, prop. 2, it suffices to consider the case where \(|\alpha|=1\) . The result then follows from the computation

\[
\begin{array}{c} \langle \partial^ {\alpha} (f g), \varphi \rangle = \langle f g, - \partial^ {\alpha} \varphi \rangle = \langle f, - g \partial^ {\alpha} \varphi \rangle \\ = \langle f, - \partial^ {\alpha} (g \varphi) + \varphi \partial^ {\alpha} g \rangle = \langle g \partial^ {\alpha} f + f \partial^ {\alpha} g, \varphi \rangle \end{array}
\]

valid for \(\varphi\in\mathcal{D}(\mathrm{U})\) .

\subsection{Schwartz functions}

We denote \(\mathcal{S}(\mathbf{R}^{n})\) the space of infinitely differentiable functions \(\varphi\) on \(\mathbf{R}^{n}\) with complex values, such that, for all \(\alpha\) and \(\beta\) in \(\mathbf{N}^{n}\) , the function \(\mathrm{X}^{\beta}\partial^{\alpha}\varphi\) is bounded on \(\mathbf{R}^{n}\) . We equip \(\mathcal{S}(\mathbf{R}^{n})\) with the locally convex topology defined by the countable family of semi-norms \((q_{\alpha,\beta})_{(\alpha,\beta)\in\mathbf{N}^{n}\times\mathbf{N}^{n}}\) , where \(q_{\alpha,\beta}\) is defined by

\[
q _ {\alpha , \beta} (\varphi) = \sup _ {x \in \mathbf {R} ^ {n}} | x ^ {\beta} \partial^ {\alpha} \varphi (x) | = \| \mathrm{X} ^ {\beta} \partial^ {\alpha} \varphi \| _ {\infty}
\]

for \(\varphi\in\mathcal{S}(\mathbf{R}^{n})\) . This topology is Hausdorff. It is also defined by the semi-norms \(\widetilde{q}_{\alpha,k}\) defined by

\[
\widetilde {q} _ {\alpha , k} (\varphi) = \sup _ {x \in \mathbf {R} ^ {n}} \| x \| ^ {k} | (\partial^ {\alpha} \varphi) (x) |.
\]

for all \(\varphi\in\mathcal{S}(\mathbf{R}^{n})\) , where \(k\in\mathbf{N}\) and \(\alpha\in\mathbf{N}^{n}\) .

We say that \(\mathcal{S}(\mathbf{R}^{n})\) is the Schwartz space or the space of Schwartz functions on \(\mathbf{R}^{n}\) .

Note. --- Let \(\varphi\in\mathcal{S}(\mathbf{R}^{n})\) . For every \(k\in\mathbf{N}\) , we have

\[
\lim _ {\| x \| \to + \infty} \| x \| ^ {k} \varphi (x) = 0,
\]

since the function \(x \mapsto \|x\|^{k+1} \varphi(x)\) is bounded.

Example. --- The function \(\gamma_{n}\) defined on \(\mathbf{R}^{n}\) by \(\gamma_{n}(x)=\exp(-\|x\|^{2})\) belongs to \(\mathcal{S}(\mathbf{R}^{n})\) . Indeed, one proves by induction on k that, for every integer \(k\in \mathbf{N}\) , there exists a polynomial \(P_{k}\in\mathbf{R}[X]\) such that \(\partial_{k}\gamma_{1}=P_{k}\gamma_{1}\) .

For all \(\alpha = (\alpha_{i})\in \mathbf{N}^{n}\) and \(\beta = (\beta_{i})\in \mathbf{N}^{n}\) , and every \(x = (x_{i})\in \mathbf{R}^{n}\) , it then follows that

\[
| (\mathrm{X} ^ {\beta} \partial^ {\alpha} \gamma_ {n}) (x) | = \prod_ {i = 1} ^ {n} | x _ {i} | ^ {\beta_ {i}} | \mathrm{P} _ {\alpha_ {i}} (x _ {i}) | \gamma_ {1} (x _ {i}),
\]

which is a bounded quantity as x varies in \(\mathbf{R}^{n}\) .

Let \(\alpha \in \mathbf{N}^n\) and \(\beta \in \mathbf{N}^n\) . If \(\varphi \in \mathcal{S}(\mathbf{R}^n)\) , then \(\mathrm{X}^{\beta}\partial^{\alpha}(\varphi)\) is still a Schwartz function; the map \(\varphi \mapsto \mathrm{X}^{\beta}\partial^{\alpha}\varphi\) thus defined from \(\mathcal{S}(\mathbf{R}^n)\) into itself is continuous.

The space \(\mathcal{S}(\mathbf{R}^{n})\) is a topological algebra. More precisely, if \(\varphi_{1}\) and \(\varphi_{2}\) belong to \(\mathcal{S}(\mathbf{R}^{n})\) , then \(\varphi_{1}\varphi_{2}\) is a Schwartz function such that

\[
q _ {\alpha , \beta} (\varphi_ {1} \varphi_ {2}) \leqslant \sum_ {0 \leqslant \gamma \leqslant \alpha} \binom{\alpha}{\gamma} q _ {\gamma , \beta} (\varphi_ {1}) q _ {\alpha - \gamma , 0} (\varphi_ {2})\tag{2}
\]

for all \(\alpha\) and \(\beta \in \mathbf{N}^n\) (prop. 1 of IV, p. 196).

The canonical inclusion of \(\mathcal{S}(\mathbf{R}^{n})\) into the space \(\mathcal{C}^{\infty}(\mathbf{R}^{n})\) , equipped with the topology described in §3 of IV, p. 200, is continuous, since

\[
\sup _ {x \in K} | \partial^ {\alpha} \varphi (x) | \leqslant q _ {\alpha , 0} (\varphi)
\]

for every compact subset K of \(\mathbf{R}^{n}\) , all \(\alpha \in \mathbf{N}^{n}\) and every Schwartz function \(\varphi\) .

Lemma 7. --- Let \(k \in \mathbf{N}\) and \(\alpha \in \mathbf{N}^n\) . For every function \(\varphi \in \mathcal{S}(\mathbf{R}^n)\) and every real number \(T > 0\), we have

\[
\widetilde {q} _ {\alpha , k} (\varphi) \leqslant \mathrm{T} ^ {k} \sup _ {\| x \| \leqslant \mathrm{T}} | \partial^ {\alpha} \varphi (x) | + \frac {1}{\mathrm{T}} \widetilde {q} _ {\alpha , k + 1} (\varphi).\tag{3}
\]

Indeed, one has

\[
\begin{array}{c} \widetilde {q} _ {\alpha , k} (\varphi) \leqslant \sup _ {\| x \| \leqslant T} \| x \| ^ {k} | \partial^ {\alpha} \varphi (x) | + \sup _ {\| x \| > T} \| x \| ^ {k} | \partial^ {\alpha} \varphi (x) | \\ \leqslant T ^ {k} \sup _ {\| x \| \leqslant T} | \partial^ {\alpha} \varphi (x) | + \frac {1}{T} \sup _ {\| x \| > T} \| x \| ^ {k + 1} | \partial^ {\alpha} \varphi (x) |. \end{array}
\]

PROPOSITION 10. --- Let B be a bounded subset of \(\mathcal{S}(\mathbf{R}^{n})\) . The topology induced on B by \(\mathcal{S}(\mathbf{R}^{n})\) coincides with the topology induced by \(\mathcal{C}^{\infty}(\mathbf{R}^{n})\) .

Since the inclusion of \(\mathcal{S}(\mathbf{R}^{n})\) into \(\mathcal{C}^{\infty}(\mathbf{R}^{n})\) is continuous, it suffices to show that, for every open subset V of \(\mathcal{S}(\mathbf{R}^{n})\) , the intersection \(V \cap B\) is open in B for the topology induced by \(\mathcal{C}^{\infty}(\mathbf{R}^{n})\) .

Let V be an open subset of \(\mathcal{S}(\mathbf{R}^{n})\) . Let \(\varphi_{0} \in \mathrm{V} \cap \mathrm{B}\) . There exist a finite set I, a family \((\alpha_{i}, k_{i})_{i \in \mathrm{I}} \in (\mathbf{N}^{n} \times \mathbf{N})^{\mathrm{I}}\) and a real number \(\varepsilon > 0\) such that V contains the set of \(\varphi \in \mathcal{S}(\mathbf{R}^{n})\) satisfying

\[
\sup _ {i \in I} \widetilde {q} _ {\alpha_ {i}, k _ {i}} (\varphi - \varphi_ {0}) \leqslant \varepsilon .
\]

Since B is bounded in \(\mathcal{S}(\mathbf{R}^{n})\) , there exists \(M > 0\) such that the seminorms \(\widetilde{q}_{\alpha_{i},k_{i} + 1}\) for \(i\in I\) are bounded by \(M\) on \(B\) . Let \(\delta >0\) and \(T > 0\) be real numbers. By inequality (3), as soon as \(\varphi \in B\) satisfies the bound

\[
\sup _ {i \in I} \sup _ {\| x \| \leqslant T} | \partial^ {\alpha_ {i}} (\varphi - \varphi_ {0}) | \leqslant \delta ,\tag{4}
\]

we have

\[
\sup _ {i \in I} \widetilde {q} _ {\alpha_ {i}, k _ {i}} (\varphi - \varphi_ {0}) \leqslant \delta T ^ {k} + \frac {2 M}{T}.
\]

Set \(T = \frac{4M}{\varepsilon}\) , then \(\delta = \frac{\varepsilon}{2T^{k}}\) . One observes that \(V \cap B\) contains the neighborhood of \(\varphi_{0}\) in B for the topology induced by \(\mathcal{C}^{\infty}(\mathbf{R}^{n})\) which is defined by (4). This concludes the proof.

COROLLARY. --- Let \((\varphi_{m})_{m\in\mathbf{N}}\) be a bounded sequence in \(\mathcal{S}(\mathbf{R}^{n})\) . For every function \(\varphi\in\mathcal{S}(\mathbf{R}^{n})\) , the following assertions are equivalent:

\begin{enumerate}
\def\labelenumi{\alph{enumi})}
\item
  The sequence \((\varphi_{m})\) converges to \(\varphi\) in \(\mathcal{S}(\mathbf{R}^n)\) ;
\item
  The sequence \((\varphi_{m})\) converges to \(\varphi\) in \(\mathcal{C}^{\infty}(\mathbf{R}^{n})\) .
\end{enumerate}

Remark. --- A sequence \((\varphi_{m})\) in \(\mathcal{C}^{\infty}(\mathbf{R}^{n})\) converges if and only if, for every \(\alpha\in \mathbf{N}^{n}\) the sequence \((\partial^{\alpha}\varphi_{m})\) converges to a function \(\varphi^{(\alpha)}\) in \(\mathcal{C}(\mathbf{R}^{n})\) endowed with the topology of compact convergence. We then have \(\varphi^{(\alpha)}=\partial^{\alpha}\varphi\) and \((\varphi_{m})\) converges to \(\varphi^{(0)}\) .

Indeed, the condition is necessary. Conversely, if the sequences \((\partial^{\alpha}\varphi_{m})\) converge to functions \(\varphi^{(\alpha)}\) for all \(\alpha\in \mathbf{N}^{n}\) , it then follows from FVR, II, p. 2, th. 1, that \(\varphi^{(\alpha)}=\partial^{\alpha}\varphi^{(0)}\) , which means that the sequence \((\varphi_{m})\) converges to \(\varphi^{(0)}\) in \(\mathcal{C}^{\infty}(\mathbf{R}^{n})\) .

PROPOSITION 11. --- The space \(\mathcal{S}(\mathbf{R}^{n})\) is a Fréchet space and a Montel space.

As the space \(\mathcal{C}^{\infty}(\mathbf{R}^{n})\) is complete (EVT, III, p. 9, example \(b\) ), the corollary of Proposition 10 implies that every Cauchy sequence in \(\mathcal{S}(\mathbf{R}^{n})\) converges in \(\mathcal{S}(\mathbf{R}^{n})\) since it is bounded and converges in \(\mathcal{C}^{\infty}(\mathbf{R}^{n})\) .

The space \(\mathcal{S}(\mathbf{R}^n)\) is therefore a Fréchet space; in particular, it is barrelled (EVT, III, p. 25, cor. of prop. 2). Let B be a bounded subset of \(\mathcal{S}(\mathbf{R}^n)\) and \((\varphi_m)_{m \in \mathbf{N}}\) a sequence with values in B. Since \(\mathcal{C}^\infty(\mathbf{R}^n)\) is a Montel space (EVT, IV, p. 18, example (4)), there exists a subsequence of \((\varphi_m)_{m \in \mathbf{N}}\) which converges in \(\mathcal{C}^\infty(\mathbf{R}^n)\) , hence in \(\mathcal{S}(\mathbf{R}^n)\) (proposition 10). Hence B is relatively compact in \(\mathcal{S}(\mathbf{R}^n)\) . It follows that \(\mathcal{S}(\mathbf{R}^n)\) is a Montel space.

\subsection{Inclusions of function spaces in the space of Schwartz functions}

PROPOSITION 12. — The space \(\mathcal{D}(\mathbf{R}^{n})\) is contained in \(\mathcal{S}(\mathbf{R}^{n})\) , and the inclusion of \(\mathcal{D}(\mathbf{R}^{n})\) into \(\mathcal{S}(\mathbf{R}^{n})\) is continuous with dense image.

Let \(B \subset \mathcal{D}(\mathbf{R}^{n})\) be a bounded subset, and let K be a compact subset of \(\mathbf{R}^{n}\) such that \(\mathrm{B} \subset \mathcal{C}_{\mathrm{K}}^{\infty}(\mathbf{R}^{n})\) . Let \(\alpha \in \mathbf{N}^{n}\) and \(k \in \mathbf{N}\) . For every function \(\varphi \in B\) , it follows that

\[
\widetilde {q} _ {\alpha , k} (\varphi) \leqslant \Bigl (\sup _ {x \in \mathrm{K}} \| x \| ^ {k} \Bigr) p _ {\alpha , \mathrm{K}} (\varphi),
\]

so B is bounded in \(\mathcal{S}(\mathbf{R}^{n})\) . The continuity of the inclusion then follows from the fact that the spaces \(\mathcal{S}(\mathbf{R}^{n})\) and \(\mathcal{D}(\mathbf{R}^{n})\) are bornological.

Let us prove that \(\mathcal{D}(\mathbf{R}^n)\) is dense in \(\mathcal{S}(\mathbf{R}^n)\) . Let B be the unit ball of \(\mathbf{R}^n\) and \(\eta \in \mathcal{D}(\mathbf{R}^n)\) a test function with support contained in 2B such that \(0 \leqslant \eta \leqslant 1\) and \(\eta(x) = 1\) for all \(x \in \mathrm{B}\) (lemma 1, \(a\) ) of IV, p. 196).

Let \(\varphi \in \mathcal{S}(\mathbf{R}^n)\) . For every integer \(m \geqslant 1\) and every \(x \in \mathbf{R}^n\) , set \(\eta_m(x) = \eta(x/m)\) . Finally define \(\varphi_m = \eta_m\varphi\) ; we have \(\varphi_m \in \mathcal{D}(\mathbf{R}^n)\) .

Since \(\partial^{\alpha}\eta_{m} = m^{-|\alpha|}(\partial^{\alpha}\eta)(x / m)\) for all \(\alpha \in \mathbf{N}^{n}\) and \(x\in \mathbf{R}^n\) , from formula (2) of IV, p. 210 we deduce that the sequence \((\varphi_{m})\) is bounded in \(\mathcal{S}(\mathbf{R}^n)\) .

Let C be a compact subset of \(R^{n}\) . The sequence \((\varphi_{m})_{m\geqslant1}\) converges to \(\varphi\) in \(\mathcal{C}_{\mathrm{K}}^{\infty}(\mathbf{R}^{n})\) since \(\varphi_{m}\) coincides with \(\varphi\) on C for all sufficiently large m. Thus the sequence \((\varphi_{m})\) converges to \(\varphi\) in \(\mathcal{C}^{\infty}(\mathbf{R}^{n})\) , and the corollary of Proposition 10 of IV, p. 211 allows us to conclude that the sequence \((\varphi_{m})\) converges to \(\varphi\) in \(\mathcal{S}(\mathbf{R}^{n})\) .

Lemma 8. --- Let B be a bounded subset of \(\mathcal{D}(\mathbf{R}^{n})\) . The topology induced on B by the topology of \(\mathcal{S}(\mathbf{R}^{n})\) coincides with the topology induced by \(\mathcal{D}(\mathbf{R}^{n})\) .

Since the inclusion of \(\mathcal{D}(\mathbf{R}^{n})\) into \(\mathcal{S}(\mathbf{R}^{n})\) is continuous, the topology on B induced by \(\mathcal{D}(\mathbf{R}^{n})\) is finer than the one induced by \(\mathcal{S}(\mathbf{R}^{n})\) . Moreover, there exists a compact subset K of \(\mathbf{R}^{n}\) such that \(\mathrm{B} \subset \mathcal{C}_{\mathrm{K}}^{\infty}(\mathbf{R}^{n})\) . For every \(\alpha \in \mathbf{N}^{n}\) , we then have \(p_{\alpha,K}(\varphi) \leqslant \widetilde{q}_{\alpha,0}(\varphi)\) , which implies that the topology induced by \(\mathcal{S}(\mathbf{R}^{n})\) is finer than the one induced by \(\mathcal{D}(\mathbf{R}^{n})\) .

PROPOSITION 13. --- Let \(p \in [1, +\infty]\) . The space \(\mathcal{S}(\mathbf{R}^{n})\) is contained in \(\mathcal{L}^{p}(\mathbf{R}^{n})\) and the canonical injection of \(\mathcal{S}(\mathbf{R}^{n})\) into \(\mathcal{L}^{p}(\mathbf{R}^{n})\) is continuous. The image of \(\mathcal{S}(\mathbf{R}^{n})\) in \(\mathrm{L}^{p}(\mathbf{R}^{n})\) is dense if \(p \neq +\infty\) .

The first assertion is immediate for \(p = +\infty\) . Suppose now that \(p \in [1, +\infty[\) . Let m be an integer such that \(n + 1 < mp\) . For every \(\varphi \in \mathcal{S}(\mathbf{R}^{n})\) and \(x \in \mathbf{R}^{n}\) , we have

\[
\| x \| ^ {n + 1} | \varphi (x) | ^ {p} \leqslant \widetilde {q} _ {0, m} (\varphi) ^ {p}
\]

hence \(\varphi \in \mathcal{L}^p (\mathbf{R}^n)\) by prop. 3 of IV, p. 199. Moreover, we obtain

\[
\mathrm{N} _ {p} (\varphi) \leqslant a _ {n} ^ {1 / p} \widetilde {q} _ {0, 0} (\varphi) + b _ {n} \widetilde {q} _ {0, m} (\varphi),
\]

where

\[
a _ {n} = \int_ {\| x \| \leqslant 1} d \mu (x), \quad b _ {n} = \int_ {\| x \| \geqslant 1} \frac {1}{\| x \| ^ {n + 1}} d \mu (x),
\]

therefore the injection of \(\mathcal{S}(\mathbf{R}^n)\) into \(\mathcal{L}^p (\mathbf{R}^n)\) is continuous.

As the space \(\mathcal{D}(\mathbf{R}^n)\) is contained in \(\mathcal{S}(\mathbf{R}^n)\) , Proposition 4 of IV, p. 202 implies that \(\mathcal{S}(\mathbf{R}^n)\) is dense in \(\mathrm{L}^p (\mathbf{R}^n)\) if \(p\neq +\infty\) .

\subsection{Functions of polynomial growth}

Definition 3. --- A function \(f\colon\mathbf{R}^{n}\to\mathbf{C}\) has polynomial growth if there exists an integer \(k\geqslant1\) such that the map defined by \(x\mapsto(1+\|x\|)^{-k}f(x)\) is bounded on \(\mathbf{R}^{n}\) .

Every function with polynomial growth is locally bounded. Every polynomial function on \(\mathbf{R}^{n}\) has polynomial growth.

PROPOSITION 14. --- Let \(f \in \mathcal{C}^{\infty}(\mathbf{R}^{n})\) . Assume that for every \(\alpha\) in \(\mathbf{N}^{n}\) , the function \(\partial^{\alpha}f\) is of polynomial growth. The linear map from the space \(\mathcal{S}(\mathbf{R}^{n})\) into itself defined by \(\varphi \mapsto f\varphi\) is continuous.

For every \(\varphi\) in \(\mathcal{S}(\mathbf{R}^{n})\) , the function \(f\varphi\) belongs to \(\mathcal{C}^{\infty}(\mathbf{R}^{n})\) . By hypothesis, for every \(\alpha \in \mathbf{N}^{n}\) , there exists an integer \(k_{\alpha} \geqslant 0\) and a real number \(\mathrm{C}_{\alpha}\) such that \(|\partial^{\alpha}f(x)| \leqslant \mathrm{C}_{\alpha}(1 + \|x\|)^{k_{\alpha}}\) for all \(x\) in \(\mathbf{R}^{n}\) . Let \(\varphi \in \mathcal{S}(\mathbf{R}^{n})\) . Let \(\alpha \in \mathbf{N}^{n}\) and \(k \in \mathbf{N}\) . By prop. 1 of IV, p. 196, it follows

\[
\begin{array}{l} \widetilde {q} _ {\alpha , k} (f \varphi) \leqslant \sum_ {0 \leqslant \beta \leqslant \alpha} \binom {\alpha} {\beta} \sup _ {x \in \mathbf {R} ^ {n}} \| x \| ^ {k} | \partial^ {\beta} f (x)   \partial^ {\alpha - \beta} \varphi (x) | \\ \leqslant \sum_ {0 \leqslant \beta \leqslant \alpha} \binom {\alpha} {\beta} C _ {\beta} \sup _ {x \in \mathbf {R} ^ {n}} \| x \| ^ {k} (1 + \| x \|) ^ {k _ {\beta}} | \partial^ {\alpha - \beta} \varphi (x) |. \end{array}
\]

Let \(\beta \in \mathbf{N}^n\) such that \(0 \leqslant \beta \leqslant \alpha\) . For every \(x \in \mathbf{R}^n\) , we have

\[
\begin{array}{c} \| x \| ^ {k} (1 + \| x \|) ^ {k _ {\beta}} | \partial^ {\alpha - \beta} \varphi (x) | \leqslant \sup _ {\| x \| \leqslant 1} 2 ^ {k _ {\beta}} | \partial^ {\alpha - \beta} \varphi (x) | \\ + \sup _ {x \in \mathbf {R} ^ {n}} 2 ^ {k _ {\beta}} \| x \| ^ {k + k _ {\beta}} | \partial^ {\alpha - \beta} \varphi (x) | \end{array}
\]

whence finally

\[
\widetilde {q} _ {\alpha , k} (f \varphi) \leqslant \sum_ {0 \leqslant \beta \leqslant \alpha} \binom {\alpha} {\beta} 2 ^ {k _ {\beta}} \mathrm{C} _ {\beta} \Bigl (\widetilde {q} _ {\alpha - \beta , 0} (\varphi) + \widetilde {q} _ {\alpha - \beta , k + k _ {\beta}} (\varphi) \Bigr),
\]

which implies the proposition.

\subsection{Tempered distributions}

DEFINITION 4. --- We call the space of tempered distributions on \(\mathbf{R}^{n}\) the dual space of \(\mathcal{S}(\mathbf{R})\) endowed with the topology of bounded convergence. It is denoted \(\mathcal{S}'(\mathbf{R}^{n})\) .

Since \(\mathcal{S}(\mathbf{R}^n)\) is bornological, the space \(\mathcal{S}'(\mathbf{R}^n)\) is complete (TVS, III, p. 24, cor. 1). Since \(\mathcal{S}(\mathbf{R}^n)\) is a Montel space, the same holds for \(\mathcal{S}'(\mathbf{R}^n)\) (EVT, IV, p. 19, prop. 9). The space \(\mathcal{S}'(\mathbf{R}^n)\) is therefore reflexive (EVT, IV, p. 16, th. 2).

Lemma 9. --- A linear map f from \(\mathcal{S}(\mathbf{R}^{n})\) into \(\mathbf{C}\) is a tempered distribution if and only if for every family \((\mathrm{M}_{\alpha,k})_{(\alpha,k)\in\mathbf{N}^{n}\times\mathbf{N}}\) in \(\mathbf{R}_{+}\) the linear form f is bounded on the set of functions \(\varphi\in\mathcal{S}(\mathbf{R}^{n})\) such that \(\widetilde{q}_{\alpha,k}(\varphi)\leqslant\mathrm{M}_{\alpha,k}\) for all \((\alpha,k)\in\mathbf{N}^{n}\times\mathbf{N}\) .

Indeed, the space \(\mathcal{S}(\mathbf{R}^{n})\) is bornological (EVT, III, p. 12, prop. 2) and every bounded subset of \(\mathcal{S}(\mathbf{R}^{n})\) is contained in one of the bounded sets described in the statement.

Lemma 10. --- Let F be a filter on \(\mathcal{S}'(\mathbf{R}^n)\) having a countable base or containing a simply bounded set. Then F converges to a tempered distribution if and only if \(\langle f, \varphi \rangle\) converges according to F for every Schwartz function \(\varphi\) .

In particular, a sequence \((f_{m})_{m\in\mathbf{N}}\) of tempered distributions converges to a tempered distribution f if and only if one has \(\langle f_{m},\varphi\rangle\to\langle f,\varphi\rangle\) for all \(\varphi\in\mathcal{S}(\mathbf{R}^{n})\) .

Since \(\mathcal{S}'(\mathbf{R}^n)\) is a Fréchet space, hence barrelled, (EVT, III, p. 25, cor. of prop. 2) the lemma follows from the Banach-Steinhaus theorem (EVT, III, p. 26, cor. 3 of th. 1).

The canonical injection \(j\) of \(\mathcal{D}(\mathbf{R}^{n})\) into \(\mathcal{S}(\mathbf{R}^{n})\) is continuous and its image is dense (Lemma 12 of IV, p. 212), hence the transpose of \(j\) , which is the restriction map of tempered distributions to the subspace \(\mathcal{D}(\mathbf{R}^{n})\) is an injective continuous linear map from \(\mathcal{S}'(\mathbf{R}^{n})\) into \(\mathcal{D}'(\mathbf{R}^{n})\) . We will identify \(\mathcal{S}'(\mathbf{R}^{n})\) with a subspace of \(\mathcal{D}'(\mathbf{R}^{n})\) via this map.

Let \(\alpha \in \mathbf{N}^n\) . The linear map \(\varphi \mapsto \partial^\alpha \varphi\) from \(\mathcal{S}(\mathbf{R}^n)\) into \(\mathcal{S}(\mathbf{R}^n)\) is continuous. Its transpose is therefore a continuous linear map from \(\mathcal{S}'(\mathbf{R}^n)\) into \(\mathcal{S}'(\mathbf{R}^n)\) (EVT, IV, p. 6, cor., \(b\) )). We denote \(\partial^\alpha\) the continuous linear map \((-1)^{|\alpha|}{}^{t}\partial^\alpha\) from \(\mathcal{S}'(\mathbf{R}^n)\) into \(\mathcal{S}'(\mathbf{R}^n)\) . This definition is compatible with definition 2 of IV, p. 208 for distributions.

Let \(h \in \mathcal{C}^{\infty}(\mathbf{R}^{n})\) be a function such that \(\partial^{\alpha}h\) is of polynomial growth for every \(\alpha \in \mathbf{N}^{n}\) . The transpose of the continuous linear map \(\varphi \mapsto h\varphi\) (prop. 14 of IV, p. 214) is a continuous linear map on \(\mathcal{S}'(\mathbf{R}^{n})\) , denoted \(f \mapsto hf\) .